%% toc: P15 169 169.1,169.2,169.3,169.4
\chapter{Likelihood Estimation under Stated Regularity}
\label{sec:p15-w02b-p15-057}

This chapter supplies the general likelihood-estimation method that
\cref{sec:p15-fisher-estimation} and \cref{sec:p15-fisher-grouping}
deliberately stop short of. Its inputs are a dominated family of sampling
laws, an independent identically distributed sampling scheme and a list of
named regularity hypotheses. Its conclusions are a conditional-law
sufficiency criterion, a strong consistency theorem for maximum likelihood,
an asymptotic normality theorem for a consistent root of the likelihood
equation, a joint version of that theorem for several parameters whose
hypotheses allow a support depending on the parameter, an information bound for unbiased statistics and a definition of
efficiency that is calibrated against that bound. Each conclusion holds only
under the hypotheses written in its statement; none follows from the bare
fact that a function has been maximized. The second half reconstructs
Fisher's own information and efficiency calculations of 1922 and records
exactly which of them the general theorems cover.

\section{Model, likelihood, information and efficiency}
\label{sec:p15-w02b-p15-057-model}

\begin{hbdefinition}{def:p15-w02b-p15-057-model}{Dominated sampling model, likelihood and maximum likelihood}
Let \((\mathsf X,\mathscr X)\) be a measurable space, \(\mu\) a
\(\sigma\)-finite measure on it and \(\Theta\) a parameter set. A
\emph{dominated model} assigns to each \(\theta\in\Theta\) a probability
measure \(P_\theta\ll\mu\) and fixes a version
\(f_\theta=dP_\theta/d\mu\) with \(0\leq f_\theta<\infty\). This is an
instance of \cref{def:p15-statistical-experiment}.
\emph{Independent identically distributed sampling} at \(\theta\) means
that observations \(X_1,X_2,\ldots\) are independent with common law
\(P_\theta\). The \emph{likelihood} and \emph{log likelihood} of the first
\(n\) observations are
\[
 L_n(\theta)=\prod_{i=1}^nf_\theta(X_i),\qquad
 \ell_n(\theta)=\sum_{i=1}^n\log f_\theta(X_i)\in[-\infty,\infty),
\]
with \(\log0=-\infty\). An \emph{estimand} is a function
\(\tau:\Theta\to\mathbb R\); a \emph{statistic} is a measurable function of
the observations alone. A \emph{maximum likelihood estimate} at sample size
\(n\) is a value \(\widehat\theta_n\in\Theta\) with
\(\ell_n(\widehat\theta_n)=\sup_{\theta\in\Theta}\ell_n(\theta)\); it may
fail to exist or to be unique. The model \emph{identifies} \(\theta\) when
\(\theta\neq\theta'\) implies \(P_\theta\neq P_{\theta'}\).
\end{hbdefinition}

The likelihood depends on the chosen density versions only through
\(\mu\)-null sets, and the sampling statements below hold almost surely,
so the version is immaterial to them. A single recorded sample, by
contrast, can sit on such a null set; existence and uniqueness of a
maximizer are therefore properties of the realized function, checked
sample by sample, not of the model.

\begin{hbdefinition}{def:p15-w02b-p15-057-regular}{Scalar regularity at an interior point}
Let \(\Theta\) be an open real interval, or contain one around the point
\(\theta_0\) of interest. The model is \emph{regular at \(\theta_0\)} when:
\begin{enumerate}
\item[(R1)] the set \(A=\{x:f_\theta(x)>0\}\) is the same for every
 \(\theta\) in a neighbourhood \(N\) of \(\theta_0\);
\item[(R2)] for each \(x\in A\), \(\theta\mapsto\log f_\theta(x)\) is twice
 continuously differentiable on \(N\); write
 \(s_\theta(x)=\partial_\theta\log f_\theta(x)\);
\item[(R3)] \(\int_A\partial_\theta f_\theta\,d\mu=0\) and
 \(\int_A\partial_\theta^2f_\theta\,d\mu=0\) at \(\theta=\theta_0\), the
 integrands being \(\mu\)-integrable on \(A\);
\item[(R4)] the \emph{Fisher information}
 \(I(\theta_0)=E_{\theta_0}s_{\theta_0}(X)^2\) is finite and positive;
\item[(R5)] there are \(\delta>0\) and a measurable \(M\) with
 \(E_{\theta_0}M(X)<\infty\) and
 \(|\partial_\theta^2\log f_\theta(x)|\leq M(x)\) for \(x\in A\) and
 \(|\theta-\theta_0|<\delta\).
\end{enumerate}
\end{hbdefinition}

Condition (R3) is the statement that the identity \(\int f_\theta\,d\mu=1\)
may be differentiated twice under the integral sign. It is a hypothesis,
not a consequence of (R2); \cref{thm:p15-w02b-p15-057-typeiii} exhibits a
family in which (R1) itself fails.

\begin{hbdefinition}{def:p15-w02b-p15-057-efficiency}{Asymptotic variance and efficiency}
Under independent identically distributed sampling from a model regular at
\(\theta_0\), an estimator sequence \(T_n\) has \emph{asymptotic variance}
\(v(\theta_0)\in[0,\infty)\) when
\(\sqrt n\,(T_n-\theta_0)\Rightarrow N(0,v(\theta_0))\) under
\(P_{\theta_0}\), with \(N(0,0)\) the point mass at zero. Its
\emph{efficiency} at \(\theta_0\) is
\[
 \operatorname{eff}_{\theta_0}(T)=\frac{1}{I(\theta_0)\,v(\theta_0)}
 \in(0,\infty],
\]
the value \(\infty\) being assigned when \(v(\theta_0)=0\). The
\emph{relative efficiency} of \(T\) with respect to \(S\) is
\(v_S(\theta_0)/v_T(\theta_0)\) when both are positive. The sequence is
\emph{efficient at \(\theta_0\)} when its efficiency equals one. For a
statistic with exact finite variance at fixed \(n\), the
\emph{finite-sample information efficiency} is
\(1/(nI(\theta_0)\operatorname{Var}_{\theta_0}T)\).
\end{hbdefinition}

Efficiency so defined is a comparison of variances in a stated limit
experiment. It is defined only for sequences possessing a normal limit at
the \(\sqrt n\) scale, it depends on \(\theta_0\), and the inequality
\(\operatorname{eff}\leq1\) is not part of the definition.
\cref{thm:p15-w02b-p15-057-hodges} shows that it can fail at a point;
\cref{thm:p15-w02b-p15-057-information} proves the finite bound for
unbiased statistics.

\section{Sufficiency in a dominated family}
\label{sec:p15-w02b-p15-057-sufficiency}

\begin{hbdefinition}{def:p15-w02b-p15-057-sufficiency}{Sufficiency by a common conditional law}
Let \(\{P_\theta:\theta\in\Theta\}\) be probability measures on
\((\mathsf X,\mathscr X)\) and \(T:\mathsf X\to\mathsf T\) a measurable map
into \((\mathsf T,\mathscr T)\). The statistic \(T\) is \emph{sufficient}
when for every \(B\in\mathscr X\) there is a
\(\sigma(T)\)-measurable \(\phi_B:\mathsf X\to[0,1]\), not depending on
\(\theta\), with
\[
 P_\theta(B\cap C)=\int_C\phi_B\,dP_\theta
 \qquad(C\in\sigma(T),\ \theta\in\Theta).
\]
That is, one function is a version of \(P_\theta(B\mid\sigma(T))\) in
the sense of \cref{def:p15-general-conditional-expectation} for every
\(\theta\) simultaneously.
\end{hbdefinition}

On a finite sample space this reduces to the criterion of
\cref{def:p15-fe-properties}: there \(\phi_{\{z\}}\) on the event
\(\{T=t\}\) is the parameter-free conditional mass of \(z\). No sample
size, independence or regularity enters the definition; it is a property
of the family of laws and the map \(T\).

\begin{hblemma}{thm:p15-w02b-p15-057-halmos-savage}{proved=annex:p15-w02b-p15-057.halmos-savage}{A countable mixture dominates a dominated family}
If every \(P_\theta\) is absolutely continuous with respect to a
\(\sigma\)-finite \(\mu\), there are \(\theta_1,\theta_2,\ldots\in\Theta\)
and weights \(c_i>0\) with \(\sum_ic_i=1\) such that
\(\lambda=\sum_ic_iP_{\theta_i}\) satisfies \(P_\theta\ll\lambda\) for
every \(\theta\in\Theta\).
\end{hblemma}

\begin{hbtheorem}{thm:p15-w02b-p15-057-factorization}{proved=annex:p15-w02b-p15-057.factorization}{The factorization criterion}
Let the family be dominated by a \(\sigma\)-finite \(\mu\). The statistic
\(T\) is sufficient if and only if there are a measurable \(h\geq0\) on
\(\mathsf X\) and, for each \(\theta\), a \(\mathscr T\)-measurable
\(g_\theta\geq0\) such that
\[
 \frac{dP_\theta}{d\mu}=g_\theta(T)\,h\qquad\mu\text{-almost everywhere}.
\]
\end{hbtheorem}

The theorem converts the conditional-law definition into a density
condition, but only for a dominated family. It concerns sufficiency of a
fixed map \(T\), not maximization: the positive finite countermodel in
\cref{thm:p15-fe-failures} has a unique maximum likelihood statistic whose
density admits no such factorization.

\section{Consistency of maximum likelihood}
\label{sec:p15-w02b-p15-057-consistency}

\begin{hbtheorem}{thm:p15-w02b-p15-057-wald}{proved=annex:p15-w02b-p15-057.wald}{Strong consistency on a compact parameter space}
Assume the model of \cref{def:p15-w02b-p15-057-model} with:
\begin{enumerate}
\item[(W1)] \(\Theta\) a compact metric space;
\item[(W2)] \(\theta\mapsto f_\theta(x)\) continuous on \(\Theta\) for
 every \(x\);
\item[(W3)] \(\theta\) identified;
\item[(W4)] at the true value \(\theta_0\),
 \(E_{\theta_0}\sup_{\theta\in\Theta}\bigl(\log f_\theta(X)
 -\log f_{\theta_0}(X)\bigr)^+<\infty\).
\end{enumerate}
Let the observations be independent with law \(P_{\theta_0}\). Then a
maximum likelihood estimate exists for every \(n\) and every sample with
\(\ell_n(\theta_0)>-\infty\), an event of probability one. Outside one null
event, every choice of maximizers \(\widehat\theta_n\) satisfies
\(\widehat\theta_n\to\theta_0\).
\end{hbtheorem}

The conclusion does not require the maximizer to be measurable, because it
holds simultaneously for all maximizers outside one null event. Compactness
is used twice, for existence and for a finite cover. The next theorem
removes it, in two ways: by a compactification on which (W2) and (W4) are
checked, or by a condition on the part of the parameter space outside a
compact set.

\begin{hbtheorem}{thm:p15-w02b-p15-057-wald-noncompact}{proved=annex:p15-w02b-p15-057.wald-noncompact}{Strong consistency on a noncompact parameter space}
Assume the model of \cref{def:p15-w02b-p15-057-model} with \(\Theta\) a
separable metric space, (W2) and (W3), and observations independent with
law \(P_{\theta_0}\). Assume one of the following.
\begin{enumerate}
\item[(C)] \emph{Compactification.} \(\Theta\) is a metric subspace of a
 compact metric space \(\overline\Theta\), and there are measurable
 \(\overline f_c:\mathsf X\to[0,\infty)\), \(c\in\overline\Theta\), with
 \(\overline f_\theta=f_\theta\) for \(\theta\in\Theta\), such that
 \(c\mapsto\overline f_c(x)\) is continuous on \(\overline\Theta\) for
 every \(x\); for each \(c\in\overline\Theta\setminus\Theta\),
 \(\int\overline f_c\,d\mu\leq1\) and \(\overline f_c\neq f_{\theta_0}\) on
 a set of positive \(\mu\)-measure; and
 \(E_{\theta_0}\sup_{c\in\overline\Theta}\bigl(\log\overline f_c(X)
 -\log f_{\theta_0}(X)\bigr)^+<\infty\).
\item[(K)] \emph{Condition at infinity.} There is a compact
 \(K\subseteq\Theta\) with
 \(E_{\theta_0}\sup_{\theta\in K}\bigl(\log f_\theta(X)-\log
 f_{\theta_0}(X)\bigr)^+<\infty\), and
 \(G(x)=\sup_{\theta\in\Theta\setminus K}\bigl(\log f_\theta(x)-\log
 f_{\theta_0}(x)\bigr)\) has \(E_{\theta_0}G^+<\infty\) and
 \(E_{\theta_0}G<0\).
\end{enumerate}
Let \(\varepsilon_n\geq0\) be numbers with \(\varepsilon_n\to0\). Outside
one null event, every sequence \(\widehat\theta_n\in\Theta\) with
\(\ell_n(\widehat\theta_n)\geq\ell_n(\theta_0)-n\varepsilon_n\) for all
large \(n\) converges to \(\theta_0\); in particular every choice of
maximum likelihood estimates does. Under (K) with \(\theta_0\in K\), and
under (C) when \(\Theta\) is open in \(\overline\Theta\), a maximum
likelihood estimate exists for all large \(n\), almost surely.
\end{hbtheorem}

Condition (C) is (W1), (W2) and (W4) of
\cref{thm:p15-w02b-p15-057-wald} transferred to \(\overline\Theta\), with
identification replaced at the added points by the requirement that a
limit of densities, which may lose mass, is not the true density. Condition
(K) is the form Wald's argument takes without a compactification: the
whole region outside \(K\) must have negative expected log likelihood
ratio in one supremum, which is stronger than negativity at each of its
points. The conclusion covers approximate maximizers, so it applies when
the supremum is not attained. Neither condition is implied by (W2) and
(W3). Continuity (W2) is itself a restriction: it fails for the Type III
family with the shape free, whose likelihood is unbounded as the end point
rises to the least observation
(\cref{thm:p15-w02b-p15-057-typeiii-joint} item~5), and the theorem says
nothing there.

\begin{hbtheorem}{thm:p15-w02b-p15-057-cauchy-line}{proved=annex:p15-w02b-p15-057.cauchy-line}{Cauchy location on the whole line}
Let \(f_\theta(x)=\pi^{-1}(1+(x-\theta)^2)^{-1}\) with respect to Lebesgue
measure and \(\Theta=\mathbb R\). With
\(\overline\Theta=[-\infty,\infty]\), metric
\(d(s,t)=|\arctan s-\arctan t|\) (where \(\arctan(\pm\infty)=\pm\pi/2\)) and
\(\overline f_{\pm\infty}=0\), condition (C) holds at every
\(\theta_0\in\mathbb R\), and
\(E_{\theta_0}\sup_c(\log\overline f_c-\log f_{\theta_0})^+
=E\log(1+Z^2)<\infty\) for a standard Cauchy variable \(Z\). A maximum
likelihood estimate over \(\mathbb R\) exists for every sample, every
choice is strongly consistent, and every measurable choice satisfies
\(\sqrt n\,(\widehat\theta_n-\theta_0)\Rightarrow N(0,2)\).
\end{hbtheorem}

This removes the compact interval of \cref{def:p15-w02b-p15-057-cauchy}.
The likelihood equation still has several roots in general, so the
conclusion is for the global maximizer.

\section{Asymptotic normality, information and efficiency}
\label{sec:p15-w02b-p15-057-asymptotic}

\begin{hbtheorem}{thm:p15-w02b-p15-057-asymptotic}{proved=annex:p15-w02b-p15-057.asymptotic}{A consistent root of the likelihood equation}
Let the model be regular at an interior \(\theta_0\)
(\cref{def:p15-w02b-p15-057-regular}) and the observations independent
with law \(P_{\theta_0}\). Let \(\widehat\theta_n\) be random variables
with \(\widehat\theta_n\to\theta_0\) in probability and
\(P_{\theta_0}(\ell_n'(\widehat\theta_n)=0)\to1\). Then
\[
 \sqrt n\,(\widehat\theta_n-\theta_0)\Rightarrow
 N\bigl(0,I(\theta_0)^{-1}\bigr),\qquad
 -\frac1n\,\ell_n''(\widehat\theta_n)\to I(\theta_0)
 \text{ in probability},
\]
and \(E_{\theta_0}\partial_\theta^2\log f_{\theta_0}(X)=-I(\theta_0)\).
The sequence is efficient at \(\theta_0\) in the sense of
\cref{def:p15-w02b-p15-057-efficiency}.
\end{hbtheorem}

\begin{hbcorollary}{thm:p15-w02b-p15-057-mle}{proved=annex:p15-w02b-p15-057.mle}{Maximum likelihood on a compact interval}
Let \(\Theta=[c,d]\) with \(c<\theta_0<d\). Assume (W2)--(W4) and
regularity at \(\theta_0\). Any measurable choice of maximum likelihood
estimates satisfies the conclusion of
\cref{thm:p15-w02b-p15-057-asymptotic}.
\end{hbcorollary}

\begin{hbtheorem}{thm:p15-w02b-p15-057-information}{proved=annex:p15-w02b-p15-057.information}{The information bound for a fixed sample size}
Assume (R1) and (R2) at \(\theta\), the first identity in (R3), and
\(0<I(\theta)<\infty\). Let \(T\) be a statistic of
\(X_1,\ldots,X_n\) with \(E_\theta T^2<\infty\), and suppose
\(m(\vartheta)=E_\vartheta T\) is differentiable at \(\theta\) with
\(m'(\theta)=E_\theta[T\,\ell_n'(\theta)]\). Then
\[
 \operatorname{Var}_\theta(T)\geq\frac{m'(\theta)^2}{n\,I(\theta)}.
\]
In particular an unbiased statistic satisfying these premises has
finite-sample information efficiency at most one.
\end{hbtheorem}

The premise \(m'(\theta)=E_\theta[T\ell_n']\) is the differentiation of
\(\int T\prod f_\vartheta(x_i)\,d\mu^n\) under the integral sign. It is
stated, not assumed silently: it fails for the endpoint family of
\cref{thm:p15-w02b-p15-057-typeiii} with the endpoint treated as unknown.
The finite version is \cref{thm:p15-fe-score}.

\begin{hbtheorem}{thm:p15-w02b-p15-057-hodges}{proved=annex:p15-w02b-p15-057.hodges}{Efficiency above one at a single point}
Let \(X_i\) be independent with mean \(\theta\in\mathbb R\) and variance
one; for the normal location family \(N(\theta,1)\), \(I\equiv1\). Put
\(\overline X_n=n^{-1}\sum_{i\leq n}X_i\) and
\(T_n=\overline X_n\) when \(|\overline X_n|\geq n^{-1/4}\), \(T_n=0\)
otherwise. Then \(\sqrt n\,T_n\to0\) in probability at \(\theta=0\), and
\(\sqrt n\,(T_n-\theta)\Rightarrow N(0,1)\) at every \(\theta\neq0\).
In the normal family \(T_n\) therefore has efficiency \(\infty\) at zero
and one elsewhere.
\end{hbtheorem}

No bound of the form \(v(\theta)\geq1/I(\theta)\) holds for all
asymptotically normal estimator sequences at every parameter value. A bound
requires a restriction on the estimators, such as the unbiasedness and
differentiation premises of \cref{thm:p15-w02b-p15-057-information}.

\begin{hbmethod}{def:p15-w02b-p15-057-procedure}{The likelihood-estimation method with its side conditions}
\begin{enumerate}
\item State the sample space, dominating measure, parameter set,
 densities, estimand and sampling scheme. Check identification
 (\cref{def:p15-w02b-p15-057-model}).
\item Form \(\ell_n\) and find a global maximizer, or report that only a
 supremum exists. A root of \(\ell_n'=0\) is a candidate only.
\item For consistency, verify (W1)--(W4), or (W2), (W3) and one of (C), (K)
 of \cref{thm:p15-w02b-p15-057-wald-noncompact}, or supply a different
 argument.
\item For an approximate variance, verify (R1)--(R5) at the parameter of
 interest and that the point is interior; the variance is then
 \(1/(nI(\theta_0))\), estimated by \(-1/\ell_n''(\widehat\theta_n)\).
\item For a sufficiency claim, exhibit the factorization of
 \cref{thm:p15-w02b-p15-057-factorization}; the maximizer's being a
 function of \(T\) is not enough.
\item Report the efficiency of a competitor only for sequences with a
 \(\sqrt n\) normal limit, and the finite bound only under the premises of
 \cref{thm:p15-w02b-p15-057-information}.
\end{enumerate}
\end{hbmethod}

\begin{hbexample}{def:p15-w02b-p15-057-cauchy}{Cauchy location on a compact interval}
Take \(\mu\) Lebesgue measure on \(\mathbb R\),
\(f_\theta(x)=\pi^{-1}(1+(x-\theta)^2)^{-1}\) and \(\Theta=[c,d]\) with
\(c<\theta_0<d\). The assertions in this example are the content of
\cref{thm:p15-w02b-p15-057-cauchy-check}.
\end{hbexample}

\begin{hbtheorem}{thm:p15-w02b-p15-057-cauchy-check}{proved=annex:p15-w02b-p15-057.cauchy-check}{The Cauchy location family meets every hypothesis}
In the example \cref{def:p15-w02b-p15-057-cauchy}, (W1)--(W4) hold and
the family is regular at \(\theta_0\) with \(M\equiv2\) and
\(I(\theta_0)=1/2\). Hence every measurable choice of maximum likelihood
estimates satisfies
\(\sqrt n\,(\widehat\theta_n-\theta_0)\Rightarrow N(0,2)\), whereas the
sample mean has, for every \(n\), the Cauchy law of one observation.
\end{hbtheorem}

The likelihood equation in this family can have several roots, so
\cref{thm:p15-w02b-p15-057-mle} is applied to the global maximizer, not to
an arbitrary root. The comparison with the sample mean repeats the failure
of \cref{thm:p15-fe-cauchy} at the level of efficiency: the mean has no
\(\sqrt n\) normal limit and so no efficiency in the sense defined here.

\section{Fisher's curvature argument and its regular reconstruction}
\label{sec:p15-w02b-p15-057-fisher-curvature}

Fisher's argument writes the sampling law of an optimum statistic
\(\theta_1\) as normal with mean \(\theta\) and standard deviation
\(\sigma\), disregarding the dependence of \(\sigma\) on \(\theta\). He
expands \(\log f\) about \(\theta_1\) with coefficients \(a\) and \(b\),
sets \(S(a)=0\) for optimum statistics, treats \(S(b)\) as \(n\bar b\) up
to a term of order \(\sqrt n\) and \(\theta-\theta_1\) as of order
\(n^{-1/2}\), and reads off \(-1/\sigma_{\theta_1}^2=n\bar b\)
\cite[pp.~328--329]{Fisher1922Foundations}. He then states the rule
\(\partial^2L/\partial\theta^2=-1/\sigma_{\hat\theta}^2\), applicable
only where \(\hat\theta\) is normally distributed, and calls the value so
found the least possible standard deviation for a statistic estimating the
same parameter \cite[pp.~331--332]{Fisher1922Foundations}.

\cref{thm:p15-w02b-p15-057-asymptotic} is the regular reconstruction. The
normal sampling law is a conclusion there, not a premise. The bound
\(|S(b)-n\bar b|=O(\sqrt n)\) is replaced by the strong law for
\(\partial_\theta^2\log f\) together with the local domination (R5).
Fisher's \(\bar b\) equals \(-I(\theta_0)\) only through the
second-derivative identity in (R3). His \(\theta-\theta_1\) of order
\(n^{-1/2}\) is the consistency premise combined with the central limit
theorem for the score. The least-variance clause is correct for unbiased
statistics under the premises of \cref{thm:p15-w02b-p15-057-information}
at each fixed \(n\). Read as a statement about every asymptotically normal
statistic, it is false at a point by \cref{thm:p15-w02b-p15-057-hodges}.

For several parameters Fisher equates the second differentials of \(L\)
with the quadratic coefficients of a bivariate normal exponent, so that
\(\sigma_{\hat\theta}^2\) is the corresponding minor divided by the Hessian
determinant \cite[p.~332]{Fisher1922Foundations}. The algebraic part of
that rule is exact.

\begin{hbtheorem}{thm:p15-w02b-p15-057-minor}{proved=annex:p15-w02b-p15-057.minor}{Fisher's minor-over-determinant rule}
Let \(J\) be a real symmetric positive definite \(k\times k\) matrix and
\(J_{(jj)}\) the matrix obtained by deleting row and column \(j\). Then
\((J^{-1})_{jj}=\det J_{(jj)}/\det J\). For \(k=2\), writing
\(J^{-1}=\begin{pmatrix}\sigma_1^2&r\sigma_1\sigma_2\\
r\sigma_1\sigma_2&\sigma_2^2\end{pmatrix}\) with \(|r|<1\),
\[
 J_{11}=\frac{1}{(1-r^2)\sigma_1^2},\qquad
 J_{22}=\frac{1}{(1-r^2)\sigma_2^2},\qquad
 J_{12}=\frac{-r}{(1-r^2)\sigma_1\sigma_2}.
\]
\end{hbtheorem}

With \(J=-\partial^2L\) these are Fisher's three displays on p.~332 with
the sign convention that his \(L\)-derivatives are the negatives of the
entries of \(J\). That the inverse expected information is the
asymptotic covariance of a joint estimate is a separate analytic claim.
\cref{thm:p15-w02b-p15-057-multi} below proves it under stated local
conditions, and \cref{thm:p15-w02b-p15-057-typeiii-joint} verifies those
conditions for Fisher's Type III curve when \(p>1\).

\section{Several parameters}
\label{sec:p15-w02b-p15-057-several}

For a vector parameter the scalar regularity of
\cref{def:p15-w02b-p15-057-regular} is replaced by two conditions on the
sample log likelihood near the true value. They are stated so that a
support depending on the parameter is not excluded in advance; whether
they hold is then checked family by family. Write \(|v|\) for the
Euclidean norm on \(\mathbb R^k\), \(\|A\|\) for the operator norm of a
\(k\times k\) matrix, \(B(\theta_0,r)\) for the closed ball of radius
\(r\), \(\nabla\ell_n\) for the gradient and \(\nabla^2\ell_n\) for the
Hessian matrix of \(\ell_n\) in \(\theta\).

\begin{hbdefinition}{def:p15-w02b-p15-057-multi-conditions}{Local score and curvature conditions}
Let \(\Theta\subseteq\mathbb R^k\), let \(\theta_0\) be an interior point
of \(\Theta\), and let the observations be independent with law
\(P_{\theta_0}\) in the model of \cref{def:p15-w02b-p15-057-model}.
\begin{enumerate}
\item[(S)] There is a measurable \(s:\mathsf X\to\mathbb R^k\) with
 \(E_{\theta_0}s(X)=0\) and finite covariance matrix
 \(J=E_{\theta_0}s(X)s(X)^{\top}\), positive definite, and there are
 events of probability tending to one on which \(\ell_n\) is differentiable
 at \(\theta_0\) with \(\nabla\ell_n(\theta_0)=\sum_{i\leq n}s(X_i)\).
\item[(C)] For every \(K<\infty\) and \(\varepsilon>0\) there are events
 \(\Omega_n(K,\varepsilon)\) with probability tending to one on which
 \(\ell_n\) is finite and twice continuously differentiable on an open set
 containing \(B(\theta_0,Kn^{-1/2})\subseteq\Theta\), and
 \[
  \sup_{\theta\in B(\theta_0,Kn^{-1/2})}
  \bigl\|n^{-1}\nabla^2\ell_n(\theta)+J\bigr\|\leq\varepsilon .
 \]
\item[(C\(^{+}\))] For every \(\varepsilon>0\) there are \(\delta>0\) and
 events with probability tending to one on which the conclusion of (C)
 holds with \(B(\theta_0,\delta)\) in place of \(B(\theta_0,Kn^{-1/2})\).
\end{enumerate}
Condition (C\(^{+}\)) implies (C), because \(Kn^{-1/2}\leq\delta\) for
large \(n\).
\end{hbdefinition}

\begin{hbdefinition}{def:p15-w02b-p15-057-multi-efficiency}{Efficiency for a function of a vector parameter}
Under (S), let \(\tau:\Theta\to\mathbb R\) be differentiable at
\(\theta_0\) with gradient \(g\). The \emph{inverse-information value}
for \(\tau\) at \(\theta_0\) is \(g^{\top}J^{-1}g\). An estimator sequence
\(T_n\) with \(\sqrt n\,(T_n-\tau(\theta_0))\Rightarrow N(0,v)\), \(v>0\),
has \emph{efficiency} \(g^{\top}J^{-1}g/v\) for \(\tau\) at \(\theta_0\).
For \(k=1\) and \(\tau(\theta)=\theta\) this is
\cref{def:p15-w02b-p15-057-efficiency}.
\end{hbdefinition}

\begin{hbtheorem}{thm:p15-w02b-p15-057-multi}{proved=annex:p15-w02b-p15-057.multi}{Joint limit of a local root of the likelihood equations}
Assume (S) and (C) of \cref{def:p15-w02b-p15-057-multi-conditions}.
\begin{enumerate}
\item For every \(\eta>0\) there is \(K<\infty\) such that, for all large
 \(n\), an event of probability at least \(1-\eta\) carries the following:
 \(\ell_n\) has exactly one stationary point in the open ball of radius
 \(Kn^{-1/2}\) about \(\theta_0\), and that point is the unique maximizer
 of \(\ell_n\) over the closed ball.
\item Let \(\widehat\theta_n\) be random vectors with
 \(\sqrt n\,(\widehat\theta_n-\theta_0)\) tight and
 \(P_{\theta_0}(\nabla\ell_n(\widehat\theta_n)=0)\to1\). Then
 \[
  \sqrt n\,(\widehat\theta_n-\theta_0)
  -J^{-1}n^{-1/2}\sum_{i\leq n}s(X_i)\to0
  \quad\text{in probability};
 \]
 for every \(t\in\mathbb R^k\),
 \(\sqrt n\,t^{\top}(\widehat\theta_n-\theta_0)\Rightarrow
 N(0,t^{\top}J^{-1}t)\); for every Borel \(\tau\) differentiable at
 \(\theta_0\) with gradient \(g\),
 \(\sqrt n\,(\tau(\widehat\theta_n)-\tau(\theta_0))\Rightarrow
 N(0,g^{\top}J^{-1}g)\); and \(-n^{-1}\nabla^2\ell_n(\widehat\theta_n)\to
 J\) in probability.
\item Under (S) and (C\(^{+}\)), random vectors with
 \(\widehat\theta_n\to\theta_0\) in probability and
 \(P_{\theta_0}(\nabla\ell_n(\widehat\theta_n)=0)\to1\) have
 \(\sqrt n\,(\widehat\theta_n-\theta_0)\) tight, so item~2 applies to them.
\end{enumerate}
\end{hbtheorem}

Item~1 asserts that a root exists close to \(\theta_0\); it does not say
how to find it without knowing \(\theta_0\), and it says nothing about
roots or maximizers far from \(\theta_0\). The conclusion of item~2 is
stated for every linear combination, which is all that the variance of a
function \(\tau(\widehat\theta_n)\) requires; joint convergence in law of
the vector is not used anywhere in this chapter.

\begin{hbtheorem}{thm:p15-w02b-p15-057-multi-bound}{proved=annex:p15-w02b-p15-057.multi-bound}{The information bound for a function of several parameters}
Assume (S) with every \(\ell_n\) differentiable at \(\theta_0\) and
\(\nabla\ell_n(\theta_0)=\sum_{i\leq n}s(X_i)\) almost surely. Let \(T\) be a
statistic of \(X_1,\ldots,X_n\) with \(E_{\theta_0}T^2<\infty\), and put
\(g=E_{\theta_0}[T\,\nabla\ell_n(\theta_0)]\). Then
\[
 \operatorname{Var}_{\theta_0}(T)\geq\frac1n\,g^{\top}J^{-1}g,
\]
with equality exactly when
\(T-E_{\theta_0}T=n^{-1}g^{\top}J^{-1}\nabla\ell_n(\theta_0)\) almost surely.
If \(T\) is unbiased for \(\tau\) and \(\vartheta\mapsto E_\vartheta T\) may
be differentiated under the integral sign at \(\theta_0\), then \(g\) is
the gradient of \(\tau\) at \(\theta_0\).
\end{hbtheorem}

\begin{hbdefinition}{def:p15-w02b-p15-057-multi-regular}{Vector regularity with a parameter-free support}
Let \(\Theta\subseteq\mathbb R^k\) contain an open ball \(N\) about
\(\theta_0\). The model is \emph{regular at \(\theta_0\)} when (R1) of
\cref{def:p15-w02b-p15-057-regular} holds on \(N\); for each \(x\in A\),
\(\theta\mapsto\log f_\theta(x)\) is twice continuously differentiable on
\(N\), with score vector \(s_\theta(x)=\nabla\log f_\theta(x)\);
\(\int_A\nabla f_{\theta_0}\,d\mu=0\) and
\(\int_A\nabla^2f_{\theta_0}\,d\mu=0\), the integrands being
\(\mu\)-integrable; \(J=E_{\theta_0}s_{\theta_0}s_{\theta_0}^{\top}\) is
finite and positive definite; and there are \(\delta>0\) and a measurable
\(M\) with \(E_{\theta_0}M(X)<\infty\) and
\(\|\nabla^2\log f_\theta(x)\|\leq M(x)\) for \(x\in A\),
\(|\theta-\theta_0|<\delta\).
\end{hbdefinition}

\begin{hbcorollary}{thm:p15-w02b-p15-057-multi-regular}{proved=annex:p15-w02b-p15-057.multi-regular}{The regular case}
Under \cref{def:p15-w02b-p15-057-multi-regular}, (S) holds with
\(s=s_{\theta_0}\) and (C\(^{+}\)) holds with
\(J=-E_{\theta_0}\nabla^2\log f_{\theta_0}(X)\). Every sequence of
random vectors that converges in probability to \(\theta_0\) and solves the
likelihood equations with probability tending to one satisfies the
conclusions of item~2 of \cref{thm:p15-w02b-p15-057-multi}.
\end{hbcorollary}

\section{Consistency, efficiency and sufficiency on a normal scale}
\label{sec:p15-w02b-p15-057-criteria}

Fisher introduces his three criteria of estimation with the scale of a
normal population estimated from an ungrouped sample. He writes
\[
 \sigma_1=\frac1n\sqrt{\frac\pi2}\,S(|x-\bar x|),\qquad
 \sigma_2=\sqrt{\frac1n\,S(x-\bar x)^2},
\]
the mean error and the mean square error, observes that both give
\(\sigma\) when calculated from the whole population, and states that the
second carries 14 per cent more weight; he defines the efficiency of a
statistic as the square of the ratio of the probable error of the most
efficient statistic to its own \cite[p.~316]{Fisher1922Foundations}.
Here \(S\) is summation over the sample. The theorem below supplies the
model, the sampling scheme and the limit theorems behind those
statements.

\begin{hbtheorem}{thm:p15-w02b-p15-057-normal-scale}{proved=annex:p15-w02b-p15-057.normal-scale}{Mean error and mean square error for a normal scale}
Let \(X_1,X_2,\ldots\) be independent with the normal law
\(N(\mu,\sigma^2)\), \(\mu\in\mathbb R\) and \(\sigma>0\) both unknown,
and let \(\sigma_1,\sigma_2\) be the statistics displayed above computed
from the first \(n\) observations. Put
\(v_1=\tfrac12(\pi-2)\sigma^2\), \(v_2=\tfrac12\sigma^2\) and
\(c=\tfrac12\sigma^2\).
\begin{enumerate}
\item \emph{Consistency.} \(\sqrt{\pi/2}\,E|X-\mu|=\sigma\) and
 \((E(X-\mu)^2)^{1/2}=\sigma\), and \(\sigma_1\to\sigma\),
 \(\sigma_2\to\sigma\) almost surely.
\item \emph{Joint limit.} For all real \(s,t\),
 \[
  \sqrt n\,\bigl(s(\sigma_1-\sigma)+t(\sigma_2-\sigma)\bigr)\Rightarrow
  N\bigl(0,\;s^2v_1+2stc+t^2v_2\bigr).
 \]
 In particular \(\sqrt n(\sigma_1-\sigma)\Rightarrow N(0,v_1)\) and
 \(\sqrt n(\sigma_2-\sigma)\Rightarrow N(0,v_2)\), and the matrix
 \(\bigl(\begin{smallmatrix}v_1&c\\c&v_2\end{smallmatrix}\bigr)\) is
 positive definite.
\item \emph{Efficiency.} The model is regular in the sense of
 \cref{def:p15-w02b-p15-057-multi-regular} at every \((\mu,\sigma)\), with
 \(J=\operatorname{diag}(1/\sigma^2,2/\sigma^2)\), and the
 inverse-information value for \(\tau(\mu,\sigma)=\sigma\) is
 \(\sigma^2/2\). In the sense of
 \cref{def:p15-w02b-p15-057-multi-efficiency} the efficiency of
 \(\sigma_2\) is one and that of \(\sigma_1\) is
 \(1/(\pi-2)=0.8760\) to four decimals; the ratio of weights
 \(v_1/v_2\) is \(\pi-2=1.1416\).
\item \emph{Correlation.} With \(r=c/\sqrt{v_1v_2}\), one has
 \(r=1/\sqrt{\pi-2}\), \(r\sqrt{v_1}=\sqrt{v_2}\) and \(r^2=v_2/v_1\),
 the efficiency of \(\sigma_1\).
\item \emph{Sufficiency.} The pair \((\bar x,S(x-\bar x)^2)\) is
 sufficient for \((\mu,\sigma)\) in the sense of
 \cref{def:p15-w02b-p15-057-sufficiency}, and \(\sigma_2\) is a function
 of it.
\end{enumerate}
\end{hbtheorem}

The 14 per cent of p.~316 is the ratio \(v_1/v_2-1=\pi-3=0.1416\) of
limit variances; the theorem gives it as a statement about normal limits
of \(\sqrt n(\sigma_i-\sigma)\) with \(\mu\) estimated, and the
replacement of \(\mu\) by \(\bar x\) inside the absolute values is shown in
the proof not to change the limit. At a fixed sample size the two
statistics are biased in different amounts and the theorem says nothing
about their exact variances. Consistency in Fisher's sense, the population
value of the statistic equalling the parameter, is item~1's first clause;
the almost sure convergence is a separate property, proved here from the
strong law.

Fisher then argues that a statistic satisfying the criterion of
sufficiency also satisfies that of efficiency: if two statistics
\(\theta_1,\theta_2\) for \(\theta\) have in large samples a joint normal
law with standard deviations \(\sigma_1,\sigma_2\) proportional to
\(n^{-1/2}\) and correlation \(r\), and the law of \(\theta_2\) given
\(\theta_1\) does not involve \(\theta\), then \(r\sigma_2=\sigma_1\), so
that \(\sigma_1<\sigma_2\) and the efficiency of \(\theta_2\) is \(r^2\)
\cite[pp.~316--317]{Fisher1922Foundations}. In this argument
\(\sigma_1,\sigma_2\) are standard deviations, not the statistics above.
The exact content of the argument is a statement about one bivariate
normal law.

\begin{hbtheorem}{thm:p15-w02b-p15-057-suff-eff}{proved=annex:p15-w02b-p15-057.suff-eff}{Fisher's sufficiency-to-efficiency argument for a bivariate normal pair}
Let \(\Theta\subseteq\mathbb R\) have at least two points, and under each
\(\theta\in\Theta\) let \((U,V)\) have the bivariate normal law with means
\((\theta,\theta)\), standard deviations \(\sigma_1,\sigma_2>0\) and
correlation \(r\in(-1,1)\), the numbers \(\sigma_1,\sigma_2,r\) not
depending on \(\theta\).
\begin{enumerate}
\item Under \(\theta\) the normal law with mean
 \(\theta+(r\sigma_2/\sigma_1)(u-\theta)\) and variance
 \(\sigma_2^2(1-r^2)\) is a version of the conditional law of \(V\) given
 \(U=u\).
\item One Markov kernel is a version of the conditional law of \(V\) given
 \(U\) under every \(\theta\in\Theta\) if and only if \(r\sigma_2=\sigma_1\).
\item If \(r\sigma_2=\sigma_1\), then \(0<r<1\), \(\sigma_1<\sigma_2\) and
 \(\sigma_1^2/\sigma_2^2=r^2\).
\end{enumerate}
\end{hbtheorem}

The theorem has three premises that Fisher's ``in large samples'' leaves
implicit: the pair is exactly bivariate normal; both coordinates are
centred at the same \(\theta\); and the parameter-free conditional law is
a property of that normal law. For statistics whose joint law is only
asymptotically normal, a finite-sample sufficiency property does not by
itself give the third premise for the limit law, and the conclusion then
compares the variances of the two limit laws only, for competitors
\(\theta_2\) whose joint limit with \(\theta_1\) is of the stated form.
\Cref{thm:p15-w02b-p15-057-normal-scale}, item~4, exhibits the identity
\(r\sqrt{v_1}=\sqrt{v_2}\) in Fisher's own example, with \(\sigma_2\), a
function of the sufficient pair of item~5, in the role of \(\theta_1\) and
\(\sigma_1\) in the role of \(\theta_2\): the efficiency of the mean error
is the square of its limiting correlation with the mean square error.

\section{Fisher's Type III calculation}
\label{sec:p15-w02b-p15-057-type-iii}

Fisher's Type III curve has element
\(df=\frac{1}{a\,p!}\bigl(\frac{x-m}{a}\bigr)^pe^{-(x-m)/a}\,dx\) on
\(x>m\), with \(p!=\Gamma(p+1)\). He compares location, scale and shape
estimated by moments with the values obtained from the second derivatives
of \(L\) \cite[pp.~332--337]{Fisher1922Foundations}. Write
\(\xi=(x-m)/a\), so that \(\xi\) has the gamma law with shape \(p+1\) and
unit scale, and \(F(p)=\frac{d^2}{dp^2}\log(p!)\), the trigamma function
at \(p+1\) \cite[p.~336]{Fisher1922Foundations}.

\begin{hbtheorem}{thm:p15-w02b-p15-057-typeiii}{proved=annex:p15-w02b-p15-057.typeiii}{The Type III information matrix and the moment efficiencies}
Let \(a>0\) and \(p>-1\).
\begin{enumerate}
\item The support \(\{x>m\}\) depends on \(m\), so (R1) fails for the
 location parameter. Per observation,
 \(E[-\partial_m^2\log f]=pE[\xi^{-2}]/a^2\) equals
 \(\frac{1}{a^2(p-1)}\) for \(p>1\), \(+\infty\) for \(0<p\leq1\), zero for
 \(p=0\) and \(-\infty\) for \(-1<p<0\); and
 \(E[-\partial_m\partial_p\log f]=\frac1{ap}\) is finite exactly when
 \(p>0\).
\item For \(p>1\) the matrix of per-observation expectations of
 \(-\partial^2\log f\) in the order \((m,a,p)\) is
 \[
  J=\begin{pmatrix}\frac{1}{a^2(p-1)}&\frac1{a^2}&\frac1{ap}\\
  \frac1{a^2}&\frac{p+1}{a^2}&\frac1a\\
  \frac1{ap}&\frac1a&F(p)\end{pmatrix},
  \quad
  \det J=\frac{1}{a^4(p-1)}\Bigl(2F(p)-\frac2p+\frac1{p^2}\Bigr)>0,
 \]
 and equals the covariance matrix of the three scores. Its inverse has
 \((p,p)\) entry \(2/(2F(p)-2/p+1/p^2)\). The \((m,a)\) block, the
 relevant matrix when \(p\) is known, has inverse with diagonal entries
 \(\tfrac12a^2(p+1)(p-1)\) and \(\tfrac12a^2\).
 The \((m,m)\) entry alone, for \(a,p\) known, inverts to
 \(a^2(p-1)\).
\item With \(a,p\) known, the moment estimate
 \(m_\mu=\overline x-a(p+1)\) is unbiased with exact variance
 \(a^2(p+1)/n\); its finite-sample information efficiency against the
 \((m,m)\) entry is \((p-1)/(p+1)=1-2/(p+1)\) for \(p>1\).
\item With \(p\) known and \(m,a\) unknown, the moment scale estimate
 \(a_\mu=\sqrt{m_2/(p+1)}\), with \(m_2\) the sample central second
 moment, satisfies
 \(\sqrt n\,(a_\mu-a)\Rightarrow N\bigl(0,\frac{(p+4)a^2}{2(p+1)}\bigr)\)
 for every \(p>-1\). Against the inverse-information value \(a^2/2\) its
 ratio is \((p+1)/(p+4)=1-3/(p+4)\) for \(p>1\).
\item Every solution of the likelihood equation in \(a\), with \(m\) and
 \(p\) at any values, satisfies \(m+a(p+1)=\overline x\). The
 corresponding moment combination \(m_\mu+a_\mu(p_\mu+1)\) also equals
 \(\overline x\).
\end{enumerate}
\end{hbtheorem}

Item~1 locates the exact regularity behind Fisher's remark that the
efficiency of location vanishes at \(p=1\)
\cite[p.~333]{Fisher1922Foundations}: for \(0<p\leq1\) the location
entry of \(J\) is infinite and for \(-1<p\leq0\) it is not a positive
number, so the inverse-information comparison is undefined there, and his
statement that location is then governed by the end point
\cite[p.~337]{Fisher1922Foundations} is the relevant regime, made exact
by \cref{thm:p15-w02b-p15-057-endpoint}. Because
(R1) fails for \(m\) at every \(p\), neither
\cref{thm:p15-w02b-p15-057-asymptotic} nor
\cref{thm:p15-w02b-p15-057-multi-regular} applies; the joint limit is
obtained instead by checking (S) and (C) directly.

\begin{hbtheorem}{thm:p15-w02b-p15-057-typeiii-joint}{proved=annex:p15-w02b-p15-057.typeiii-joint}{The joint Type III limit for \(p>1\)}
Let \(\Theta=\mathbb R\times(0,\infty)\times(-1,\infty)\) with
coordinates \((m,a,p)\), let \(\theta_0=(m,a,p)\) with \(a>0\) and
\(p>1\), and let the observations be independent with the Type III law at
\(\theta_0\). Let \(J\) be the matrix of
\cref{thm:p15-w02b-p15-057-typeiii}, item~2.
\begin{enumerate}
\item Conditions (S) and (C) of
 \cref{def:p15-w02b-p15-057-multi-conditions} hold, with
 \(s=\bigl((1-p/\xi)/a,\ (\xi-p-1)/a,\ \log\xi-\psi(p+1)\bigr)\),
 \(\psi\) the derivative of \(\log\Gamma\). Hence item~1 of
 \cref{thm:p15-w02b-p15-057-multi} holds: with probability tending to
 one the likelihood equations have a root within \(Kn^{-1/2}\) of
 \(\theta_0\), unique there.
\item Every sequence of roots \(\widehat\theta_n=(\hat m,\hat a,\hat p)\)
 with \(\sqrt n\,(\widehat\theta_n-\theta_0)\) tight satisfies
 \[
  n\,\sigma_{\hat p}^2\to\frac{2}{2F(p)-2/p+1/p^2},\qquad
  n\,\sigma_{\hat m}^2\to a^2\,\frac{(p-1)\bigl((p+1)F(p)-1\bigr)}
  {2F(p)-2/p+1/p^2},
 \]
 in the sense that these are the variances of the normal limits of
 \(\sqrt n\,(\hat p-p)\) and \(\sqrt n\,(\hat m-m)\); and
 \(\sqrt n\,\bigl(\hat m+\hat a(\hat p+1)-m-a(p+1)\bigr)\Rightarrow
 N\bigl(0,a^2(p+1)\bigr)\).
\item The sample mean \(\overline x\), which equals
 \(m_\mu+a_\mu(p_\mu+1)\), is unbiased for \(m+a(p+1)\) and its exact
 variance \(a^2(p+1)/n\) equals the bound of
 \cref{thm:p15-w02b-p15-057-multi-bound} for that function. Its
 efficiency for \(m+a(p+1)\) in the sense of
 \cref{def:p15-w02b-p15-057-multi-efficiency} is one.
\item The ratio of the limit variance of \(\hat m\) in item~2 to the value
 \(a^2(p-1)\) obtained when \(a\) and \(p\) are known is
 \[
  n'=\frac{(p+1)F(p)-1}{2F(p)-2/p+1/p^2}.
 \]
\item For every sample, \(\sup_{\Theta}\ell_n=+\infty\), so no global
 maximum likelihood estimate exists on \(\Theta\).
\end{enumerate}
\end{hbtheorem}

Items~2 and~4 are Fisher's values of \(\sigma_{\hat p}^2\),
\(\sigma_{\hat m}^2\) and \(n'\) \cite[pp.~335--337]{Fisher1922Foundations},
now with their premises: \(p>1\) and a root sequence within
\(O(n^{-1/2})\) of the true value. Item~3 is his 100 per cent efficiency of
\(m_\mu+a_\mu(p_\mu+1)\) \cite[p.~337]{Fisher1922Foundations}; the
theorem shows that it holds for \(p>1\) both against the joint limit of
item~2 and, more strongly, at every fixed \(n\) against the information
bound. Item~5 corrects the description of the joint solution of the
likelihood equations as the set of values with the maximum likelihood
\cite[p.~335]{Fisher1922Foundations}: when \(p\) is free to range over
\((-1,\infty)\), letting \(m\) rise to the least observation with \(p<0\)
fixed sends the likelihood to infinity, so the solution Fisher computes is
a local root, not a global maximum. The large-\(p\) expansions on
pp.~335--337 and the 80 per cent threshold of p.~336 are approximations
to the exact expressions in items~2 and~4; they are made exact, with
remainders, in \cref{thm:p15-w02b-p15-057-typeiii-largep}.

\begin{hbtheorem}{thm:p15-w02b-p15-057-moment-shape}{proved=annex:p15-w02b-p15-057.moment-shape}{The moment estimate of the Type III shape}
Let \(a>0\) and \(p>-1\), and let \(m_2,m_3\) be the sample central second
and third moments. Fisher's moment equation \(4/(p+1)=\beta_1\)
\cite[p.~334]{Fisher1922Foundations}, with \(\beta_1\) replaced by
\(m_3^2/m_2^3\), gives \(p_\mu=4m_2^3/m_3^2-1\) when \(m_3\neq0\); give
\(p_\mu\) any fixed value when \(m_3=0\). Then
\[
 \sqrt n\,(p_\mu-p)\Rightarrow N\bigl(0,\,6(p+1)(p+2)(p+6)\bigr).
\]
For \(p>1\) the efficiency of \(p_\mu\) against the joint limit of
\cref{thm:p15-w02b-p15-057-typeiii-joint} is
\[
 \frac{2}{\bigl(2F(p)-2/p+1/p^2\bigr)\,6(p+1)(p+2)(p+6)} .
\]
\end{hbtheorem}

The limit variance is Fisher's \(\sigma_{p_\mu}^2\)
\cite[p.~335]{Fisher1922Foundations}; it is derived here directly from
the sixth-order moments of the gamma law, not from the sampling variance
of \(\beta_1\) that Fisher quotes \cite[p.~334]{Fisher1922Foundations}.
Evaluating the efficiency formula with the trigamma function gives
\(0.0871\), \(0.1532\), \(0.2159\) and \(0.2727\) at \(p=2,3,4,5\), Fisher's
table entries \cite[p.~336]{Fisher1922Foundations}, which are therefore
information ratios. His entries \(0\) and \(0.0274\) at \(p=0\) and
\(p=1\) are the limit and the value of the same formula where the location
entry of \(J\) is not finite; they are not information ratios, and the
theorem gives no comparison for \(-1<p\leq1\).

For large \(p\) Fisher replaces the trigamma expressions in
\cref{thm:p15-w02b-p15-057-typeiii-joint,thm:p15-w02b-p15-057-moment-shape}
by series in \(1/p\): \(2F(p)-2/p+1/p^2\) by
\(\tfrac13(p^{-3}-\tfrac15p^{-5}+\tfrac17p^{-7}\cdots)\), hence
\(\sigma_{\hat p}^2\) approximately \(\frac6n(p^3+\tfrac15p)\) and the
efficiency of the moment shape approximately
\(\frac{p^3+\frac15p}{(p+1)(p+2)(p+6)}\), with efficiencies over 80 per
cent when \(p\) exceeds \(38{\cdot}1\), \(\beta_1=0{\cdot}102\)
\cite[pp.~335--336]{Fisher1922Foundations}; and \((p+1)F(p)-1\) by
\(\frac1{2p}(1-\frac2{3p}+\frac1{3p^2}\cdots)\), hence
\(n'=\tfrac32p^2-p+\tfrac45\) approximately, the multiplier
\((p-\tfrac13)\sqrt{3/2}\) for the probable error of the end point, and
\(n'\) falling to unity at \(p=1\) \cite[p.~337]{Fisher1922Foundations}.
The series are asymptotic, not convergent, so each is stated below with
an explicit remainder.

\begin{hblemma}{thm:p15-w02b-p15-057-trigamma}{proved=annex:p15-w02b-p15-057.trigamma}{Series, integral and two-sided bounds for the trigamma function}
For \(x>0\),
\[
 \psi'(x)=\sum_{j=0}^{\infty}\frac1{(x+j)^2}
 =\int_0^\infty\frac{t\,e^{-xt}}{1-e^{-t}}\,dt ,
\]
and
\[
 \frac1x+\frac1{2x^2}+\frac1{6x^3}-\frac1{30x^5}<\psi'(x)<
 \frac1x+\frac1{2x^2}+\frac1{6x^3}-\frac1{30x^5}+\frac1{42x^7}.
\]
Consequently, for \(p>0\), \(F(p)=\psi'(p+1)\) satisfies
\(F(p)>1/(p+1)\) and
\[
 F(p)=\frac1p-\frac1{2p^2}+\frac1{6p^3}-\frac1{30p^5}+\frac{\vartheta}{42p^7}
 \quad\text{for some }\vartheta=\vartheta(p)\in(0,1).
\]
\end{hblemma}

\begin{hbtheorem}{thm:p15-w02b-p15-057-typeiii-largep}{proved=annex:p15-w02b-p15-057.typeiii-largep}{Fisher's large-\(p\) expansions with remainders, and the 80 per cent threshold}
Let \(p\geq1\), put \(D(p)=2F(p)-2/p+1/p^2\) and \(\mathcal N(p)=(p+1)F(p)-1\), and
let \(\vartheta\in(0,1)\) be as in \cref{thm:p15-w02b-p15-057-trigamma}.
For \(p>1\), \(2/D(p)\) is the limit variance of \(\sqrt n(\hat p-p)\) and
\(n'=\mathcal N(p)/D(p)\) the ratio of item~4 of
\cref{thm:p15-w02b-p15-057-typeiii-joint}; at \(p=1\) they are values of
the same formulas.
\begin{enumerate}
\item \(D(p)=\frac13\bigl(p^{-3}-\frac15p^{-5}+\frac{\vartheta}{7}p^{-7}\bigr)\)
 and
 \(\mathcal N(p)=\frac1{2p}\bigl(1-\frac2{3p}+\frac1{3p^2}-\frac1{15p^3}
 -\frac1{15p^4}+\frac{\vartheta(p+1)}{21p^6}\bigr)\).
\item \(\bigl|2/D(p)-6(p^3+\tfrac15p)\bigr|<6/(7p)\).
\item The efficiency \(E(p)\) of the moment shape in
 \cref{thm:p15-w02b-p15-057-moment-shape} satisfies
 \[
  \Bigl|E(p)-\frac{p^3+\frac15p}{(p+1)(p+2)(p+6)}\Bigr|
  <\frac1{7p(p+1)(p+2)(p+6)} .
 \]
\item \(E(p)<0.8\) for \(1<p\leq38.1\) and \(E(p)>0.8\) for
 \(p\geq38.11\). At \(p=38.1\) the moment equation of
 \cref{thm:p15-w02b-p15-057-moment-shape} gives
 \(\beta_1=4/(p+1)=0.1023\) to four decimals.
\item \(n'>1\) for \(p>1\), the formula gives \(n'=1\) at \(p=1\), and
 \(\tfrac32p^2-p+\tfrac45-\tfrac{11}{14p}<n'<\tfrac32p^2-p+\tfrac45\).
\item For \(p\geq2\),
 \(0<\sqrt{n'}-\sqrt{3/2}\,(p-\tfrac13)<
 \frac{19\sqrt{3/2}}{90\,(p-\frac13)}\).
\end{enumerate}
\end{hbtheorem}

Item~1 confirms each coefficient Fisher prints and bounds the first one he
omits. Item~4 makes his threshold exact in both directions: the crossing
of \(0.8\) lies between \(38.1\) and \(38.11\), so ``over 80 per cent when
\(p\) exceeds \(38{\cdot}1\)'' is correct as stated, the approximation of
item~3 being within \(6\times10^{-8}\) there. Item~5 replaces the falling
of \(n'\) to unity by the exact statement that \(n'\) exceeds one for every
\(p>1\); it does not assert that \(n'\) is monotone. Because \(n'\) is a
ratio of variances, the end point's probable error is \(\sqrt{n'}\) times
that of location for a curve of known form and size, and item~6 shows
that Fisher's multiplier \((p-\frac13)\sqrt{3/2}\) is accurate to order
\(1/p\). None of these items applies to \(p<1\), where the location entry
of \(J\) is not finite (\cref{thm:p15-w02b-p15-057-typeiii}, item~1).

\section{Location and scaling of a curve of given form}
\label{sec:p15-w02b-p15-057-location-scale}

Fisher specifies the form of an error curve by \(df\propto
e^{\phi(\xi)}d\xi\) with \(\xi=(x-m)/a\), obtains
\(\sigma_{\hat m}^2=-a^2/(n\bar{\phi''})\), calls \(-\bar{\phi''}/a^2\) the
intrinsic accuracy, and shows that the Cauchy curve has
\(\bar{\phi''}=-\tfrac12\) and so the same intrinsic accuracy as a normal
curve with \(\sigma^2=2a^2\) \cite[pp.~338--339]{Fisher1922Foundations}.
He then moves the origin so that the cross second differential has mean
zero, calling that point the centre of location
\cite[pp.~339--340]{Fisher1922Foundations}.

\begin{hbtheorem}{thm:p15-w02b-p15-057-location-scale}{proved=annex:p15-w02b-p15-057.location-scale}{Location--scale information and the centre of location}
Let \(g=e^{\phi}\) be a probability density on an open interval
\(D\subseteq\mathbb R\), and \(f_{m,a}(x)=a^{-1}g((x-m)/a)\). Assume
\(\phi\) is twice continuously differentiable on \(D\); that \(g\),
\(\xi g\), \(g'\), \(\xi g'\) and \(\xi^2g'\) tend to zero at both ends of
\(D\); that \(g'\), \(\xi g'\), \(g''\), \(\xi g''\) and \(\xi^2g''\) are
Lebesgue integrable on \(D\); and that \(E\phi'(\xi)^2\),
\(E\xi^2\phi'(\xi)^2\), \(E|\phi''(\xi)|\) and \(E\xi^2|\phi''(\xi)|\) are
finite, \(\xi\) having density \(g\). Then
\(E\phi'=0\), \(E\xi\phi'=-1\), and the per-observation information
matrix in \((m,a)\) is
\[
 \frac1{a^2}\begin{pmatrix}-E\phi''&-E\xi\phi''\\
 -E\xi\phi''&1-E\xi^2\phi''\end{pmatrix},
\]
with \(-E\phi''=E\phi'^2\). Fisher's displays
\(\partial^2L/\partial m\,\partial a=a^{-2}S(\xi\phi'')\) and
\(\partial^2L/\partial a^2=a^{-2}S(\xi^2\phi''-1)\) hold as equalities of
expectations, not sample by sample. If \(E\phi''\neq0\), replacing
\(\xi\) by \(\xi-c\) with \(c=E\xi\phi''/E\phi''\), that is, measuring
location from the point \(m+ac\), makes the off-diagonal entry zero. For
the Cauchy form \(\phi(\xi)=-\log(1+\xi^2)\) on \(D=\mathbb R\),
\(E\phi''=-\tfrac12\); for the Type III form \(\phi=p\log\xi-\xi\) on
\((0,\infty)\) with \(p>1\), \(E\phi''=-1/(p-1)\), \(E\xi\phi''=-1\),
\(c=p-1\), and after the shift \(1-E\xi^2\phi''=2\).
\end{hbtheorem}

The Cauchy value gives Fisher's \(\sigma_{\hat m}^2=2a^2/n\), and
\cref{thm:p15-w02b-p15-057-cauchy-check} shows that this value is the
actual asymptotic variance of the maximum likelihood location estimate
when \(a\) is known and the parameter set is a compact interval. The Type
III value \(1-E\xi^2\phi''=2\) gives his \(\sigma_a^2=a^2/(2n)\) from the
diagonal at the centre of location. The boundary hypotheses on \(g\) are
exactly what fails at the Type III end point when \(p\leq1\).

\section{Type IV and Type VII curves: location and scaling by moments}
\label{sec:p15-w02b-p15-057-type-iv}

Fisher applies the location--scale calculation to Pearson's Type IV form
\(df\propto e^{-\nu\tan^{-1}\xi}(1+\xi^2)^{-(r+2)/2}\,d\xi\), finds the
centre of location, the intrinsic accuracy and the sampling variances of
the moment estimates of location and scale, and specializes to the
symmetric Type VII curve at \(\nu=0\)
\cite[pp.~340--342]{Fisher1922Foundations}. His numerical illustration is
the value \(r=7\), which he describes as the conventional boundary for
heterotypic curves \cite[p.~341]{Fisher1922Foundations}.

\begin{hbtheorem}{thm:p15-w02b-p15-057-typeiv}{proved=annex:p15-w02b-p15-057.typeiv}{Type IV information and the efficiency of moments}
Let \(r>-1\) and \(\nu\in\mathbb R\), let \(g(\xi)=C\,e^{\phi(\xi)}\) on
\(\mathbb R\) with
\(\phi(\xi)=-\nu\tan^{-1}\xi-\frac{r+2}2\log(1+\xi^2)\) and \(C\) the
normalizing constant, and let \(f_{m,a}(x)=a^{-1}g((x-m)/a)\) with
\(m\in\mathbb R\), \(a>0\), the observations being independent with law
\(f_{m,a}\) and \(r,\nu\) known. Put \(D=(r+4)^2+\nu^2\).
\begin{enumerate}
\item The hypotheses of \cref{thm:p15-w02b-p15-057-location-scale} hold,
 and
 \[
  E\phi''=-\frac{(r+1)(r+2)(r+4)}{D},\qquad
  E\xi\phi''=\frac{\nu(r+1)(r+2)}{D},\qquad
  E\xi^2\phi''-1=-\frac{(r+1)\bigl(2(r+4)+\nu^2\bigr)}{D}.
 \]
 The centre of location is the point \(\xi=-\nu/(r+4)\), that is
 \(x=m-a\nu/(r+4)\); the mode is the point \(\xi=-\nu/(r+2)\).
\item The per-observation information for \(m\) with \(a\) known is
 \((r+1)(r+2)(r+4)/(a^2D)\). With \(m\) and \(a\) both unknown, the
 information matrix is positive definite and the \((a,a)\) entry of its
 inverse is \(a^2(r+4)/(2(r+1))\), which does not involve \(\nu\). The
 model in \((m,a)\) satisfies \cref{def:p15-w02b-p15-057-multi-regular} at
 every parameter value, so by \cref{thm:p15-w02b-p15-057-multi-regular}
 every consistent root of the likelihood equations has these inverse
 information values as limit variances.
\item Let \(r>1\). Then \(E\xi=-\nu/r\) and
 \(\operatorname{Var}\xi=(r^2+\nu^2)/(r^2(r-1))\). With \(a\) known, the
 moment estimate \(m_\mu=\overline x+a\nu/r\) is unbiased with variance
 \(a^2(r^2+\nu^2)/(nr^2(r-1))\), and its finite-sample information
 efficiency, which is also its efficiency in the sense of
 \cref{def:p15-w02b-p15-057-efficiency}, is
 \[
  \frac{r^2(r-1)\bigl((r+4)^2+\nu^2\bigr)}{(r+1)(r+2)(r+4)(r^2+\nu^2)} .
 \]
 At \(\nu=0\) this is \(1-6/((r+1)(r+2))\); at \(r=7\) it is
 \(\frac{49}{132}\cdot\frac{121+\nu^2}{49+\nu^2}\), which decreases in
 \(\nu^2\) from \(121/132\) at \(\nu=0\) towards \(49/132\). For
 \(0<r\leq1\) the variance of one observation is infinite, and for
 \(-1<r\leq0\) its mean does not exist.
\item Let \(r>3\), and let \(m_2\) be the sample central second moment.
 The kurtosis of one observation is
 \[
  \beta_2=\frac{3(r-1)}{(r-2)(r-3)}\Bigl(r+6-\frac{8r^2}{r^2+\nu^2}\Bigr),
 \]
 and with \(m\) unknown the moment estimate
 \(a_\mu=\sqrt{m_2/\operatorname{Var}\xi}\) is such that
 \(\sqrt n\,(a_\mu-a)\) converges in law to
 \(N\bigl(0,a^2(\beta_2-1)/4\bigr)\), where
 \[
  \frac{\beta_2-1}{4}=\frac{r^3(r-2)+\nu^2(r^2+10r-12)}
  {2(r-2)(r-3)(r^2+\nu^2)} .
 \]
 Its efficiency for \(a\) in the sense of
 \cref{def:p15-w02b-p15-057-multi-efficiency} is
 \[
  \frac{(r-2)(r-3)(r+4)(r^2+\nu^2)}
  {(r+1)\bigl(r^3(r-2)+\nu^2(r^2+10r-12)\bigr)} ;
 \]
 at \(\nu=0\) this is \(1-12/(r(r+1))\), and at \(r=7\) it is
 \(\frac{55}2\cdot\frac{49+\nu^2}{1715+107\nu^2}\), which decreases in
 \(\nu^2\) from \(2695/3430\) at \(\nu=0\) towards \(55/214\).
\end{enumerate}
\end{hbtheorem}

Items~3 and~4 are Fisher's displays (3) and (4), his Type VII
specializations and his \(r=7\) values
\cite[pp.~341--342]{Fisher1922Foundations}; \(121/132\) and \(49/132\) are
his 91.67 and 37.12 per cent, and \(2695/3430\) and \(55/214\) his 78.57
and 25.70 per cent. Two qualifications attach to them. First, the
location comparison assumes \(a\), \(r\) and \(\nu\) known and the scale
comparison assumes \(r\) and \(\nu\) known; with the form parameters also
estimated the relevant inverse-information values change. Second, the
location efficiency in item~3 is a fixed-\(n\) comparison of an unbiased
statistic with the bound of \cref{thm:p15-w02b-p15-057-information}, whereas
the scale efficiency in item~4 is a ratio of limit variances. Fisher
states that the centre of location lies at the distance \(-\nu a/(r+4)\)
from the mode \cite[p.~340]{Fisher1922Foundations}. By item~1 that
displacement is measured from the point \(\xi=0\) of the displayed form;
measured from the mode it is \(2\nu a/((r+2)(r+4))\). The two agree only
for the symmetric curve \(\nu=0\), where the centre of location, the mode
and the mean coincide.

\section{The Pearson curves: where moments are valid and how efficient they are}
\label{sec:p15-w02b-p15-057-pearson}

Fisher's Section~10 surveys the whole Pearson system, draws the regions in
which the first and second moments are usable estimates of location and
scale, approximates the efficiency contours near the normal curve, treats
curves that end abruptly through the rectangular law, and computes the
efficiency of moments for the form parameters of the symmetric Types~VII
and~II \cite[pp.~342--355]{Fisher1922Foundations}. This section supplies
the parts of that survey that can be stated and proved from the
calculations of \cref{sec:p15-w02b-p15-057-type-iv}, and records Fisher's
form-parameter results with their status.

\begin{hbdefinition}{def:p15-w02b-p15-057-pearson-system}{The Pearson system and Fisher's diagram}
A \emph{Pearson curve} is a probability density \(y\) on an interval
satisfying Fisher's equation~(5),
\(y'/y=-(x-m)/(a+bx+cx^2)\) \cite[p.~342]{Fisher1922Foundations}.
According as the quadratic \(a+bx+cx^2\) has real or complex roots,
the solutions are, after a shift of origin, of the form
\(df\propto(1+x/a_1)^{m_1}(1-x/a_2)^{m_2}\,dx\) or of the Type~IV form
\(df\propto(1+x^2/a^2)^{-(r+2)/2}e^{-\nu\tan^{-1}(x/a)}\,dx\) of
\cref{thm:p15-w02b-p15-057-typeiv}. Integrability requires an exponent
greater than \(-1\) at a finite end and, for an infinite range, decay of
the ordinate faster than \(1/x\) \cite[p.~343]{Fisher1922Foundations}.
Fisher's named special cases are Type~II (symmetric, finite range, with
\(r\) negative in his convention), Type~III (\(df\propto x^pe^{-x}dx\)),
Type~V, Type~VI and Type~VII (Type~IV with \(\nu=0\))
\cite[pp.~344--345]{Fisher1922Foundations}. In Fisher's conspectus
(figs.~2--4) a curve with parameters \((r,\nu)\) is plotted at the point
\((x,y)\) determined by
\[
 r=y-\frac1y,\qquad
 r^2+\nu^2=\frac{\bigl((1+y)^2+x^2\bigr)\bigl((1-y)^2-x^2\bigr)}{y\,(x^2+y)},
\]
so that lines of constant \(r\) are horizontal, the symmetric curves lie on
\(x=0\) and the normal curve is the point \(N=(0,0)\)
\cite[p.~345]{Fisher1922Foundations}. Solving, at a point with \(y\neq0\)
and \(x^2+y\neq0\),
\[
 \nu^2=-\frac{x^2\bigl(x^2y+(1+y^2)^2\bigr)}{y^2\,(x^2+y)},
\]
which is positive exactly when \(x\neq0\) and \(x^2y+(1+y^2)^2\) and
\(x^2+y\) have opposite signs; near \(N\) this is the region \(y<0\),
\(0<x^2<-y\), the Type~IV side of the Type~VII line.
\end{hbdefinition}

\begin{hbtheorem}{thm:p15-w02b-p15-057-pearson-moments}{proved=annex:p15-w02b-p15-057.pearson-moments}{Moments of the Type IV curve and the boundaries they impose}
Let \(\xi\) have the Type~IV density \(g\) of
\cref{thm:p15-w02b-p15-057-typeiv}, with \(r>-1\) and \(\nu\) real, and
write \(\mu_k'=E\xi^k\).
\begin{enumerate}
\item For \(k\geq0\), \(E|\xi|^k<\infty\) if and only if \(k<r+1\).
\item If \(k\geq0\) and \(k+1<r+1\), then
 \((r-k)\,\mu_{k+1}'=k\,\mu_{k-1}'-\nu\,\mu_k'\) (with \(\mu_{-1}'=0\)).
\item Hence the variance of a sample mean is finite if and only if
 \(r>1\); the central fourth moment, and so \(\beta_2\), is finite if and
 only if \(r>3\); and, for \(j\geq1\), the sample moment
 \(n^{-1}\sum_{i\leq n}\xi_i^j\) about a fixed origin has finite sampling
 variance if and only if \(2j<r+1\), so for the second moment if and only
 if \(r>3\) and for the fourth moment if and only if \(r>7\).
\item For \(r>3\),
 \[
  \beta_1=\frac{16\nu^2(r-1)}{(r-2)^2(r^2+\nu^2)},\qquad
  \frac{\beta_1(\beta_2+3)^2}{4(4\beta_2-3\beta_1)(2\beta_2-3\beta_1-6)}
  =\frac{\nu^2}{r^2+\nu^2},
 \]
 with \(\beta_2\) as in item~4 of \cref{thm:p15-w02b-p15-057-typeiv}.
\end{enumerate}
\end{hbtheorem}

Item~3 gives exact content to the lines of Fisher's figs.~2--4: the base
\(r=1\) of the region of validity of the first moment, the line \(r=3\)
below which \(\beta_2\) is infinite, and the heterotypic limit \(r=7\),
between which and \(r=3\) the curves have finite \(\beta_1,\beta_2\) while
the fourth moment of a sample has infinite sampling variance
\cite[pp.~345--347]{Fisher1922Foundations}. Item~4 is Pearson's criterion
\(\kappa_2\) and Fisher's identity \(\kappa_2=1-r^2/(4\epsilon)\) with
\(4\epsilon=r^2+\nu^2\) \cite[p.~345]{Fisher1922Foundations}. The moment
expression on the left is defined only for \(r>3\); Fisher's objection
that the criterion would put all curves with \(r=0\) into Type~V concerns
the right-hand side \(1-r^2/(4\epsilon)\) extended beyond that range, not
the moment ratio, which does not exist there.

\begin{hbtheorem}{thm:p15-w02b-p15-057-contours}{proved=annex:p15-w02b-p15-057.contours}{The efficiency contours near the normal point}
Let \(E_1(r,\nu)\) be the location efficiency of the mean in item~3 of
\cref{thm:p15-w02b-p15-057-typeiv} and \(E_2(r,\nu)\) the scale efficiency
of the second moment in item~4 there, and express both in Fisher's
coordinates \((x,y)\) of \cref{def:p15-w02b-p15-057-pearson-system}. Then
\[
 1-E_1=\frac{2\bigl(3r^3+12r^2+\nu^2(4r^2+7r+4)\bigr)}
 {(r^2+\nu^2)(r+1)(r+2)(r+4)},\qquad
 1-E_2=\frac{12\bigl(r^3-2r^2+\nu^2(r^2+r-3)\bigr)}
 {(r+1)\bigl(r^3(r-2)+\nu^2(r^2+10r-12)\bigr)},
\]
and each of \(1-E_1\), \(1-E_2\), as a function of \((x,y)\), is a ratio of
polynomials whose denominator equals \(1\) at \(N\). Consequently there
are \(\rho>0\) and \(C<\infty\) with
\[
 \bigl|1-E_1-(8x^2+6y^2)\bigr|\leq C\,(x^2+y^2)^{3/2},\qquad
 \bigl|1-E_2-12(x^2+y^2)\bigr|\leq C\,(x^2+y^2)^{3/2}
\]
for \(x^2+y^2<\rho^2\), on the part of that disc where the curves are of
Type~IV or VII (\(y<0\), \(x^2<-y\)).
\end{hbtheorem}

These are Fisher's ellipses \(8x^2+6y^2=1-E\) for the first moment and
circles \(12x^2+12y^2=1-E\) for the second
\cite[p.~347]{Fisher1922Foundations}, now with an explicit error order. On
the other side of \(N\) (\(y>0\)) the curves are of Types~I and~II and the
same rational functions are an analytic continuation with \(\nu^2<0\); that
they are the efficiencies of those curves is not shown here, so the
contours of figs.~3 and~4 above \(N\) rest on Fisher's calculation alone.

The upper boundary of the region of validity of the mean is, in Fisher's
words, formed by the curves meeting the axis at a finite angle at an end of
their range, \(m_1\) or \(m_2\) equal to \(1\)
\cite[p.~346]{Fisher1922Foundations}; he explains it by the behaviour of
observations near the terminus \cite[pp.~347--348]{Fisher1922Foundations}.

\begin{hbtheorem}{thm:p15-w02b-p15-057-terminus}{proved=annex:p15-w02b-p15-057.terminus}{Curves that end at a finite point}
\begin{enumerate}
\item Let \(F\) be a distribution function with \(F(m)=0\) and
 \(F(m+x)=k'x^{\alpha+1}(1+o(1))\) as \(x\downarrow0\), where
 \(\alpha>-1\) and \(k'>0\), and let \(X_{(1)}\) be the least of \(n\)
 independent observations with law \(F\). Then for every \(t\geq0\),
 \(P\bigl(n^{1/(\alpha+1)}(X_{(1)}-m)>t\bigr)\to\exp(-k't^{\alpha+1})\).
 If \(\alpha<1\), \(\sqrt n\,(X_{(1)}-m)\to0\) in probability.
\item Let \(f(x)=C(1+x/a_1)^{m_1}(1-x/a_2)^{m_2}\) on \((-a_1,a_2)\), with
 \(a_1,a_2>0\), \(m_1,m_2\in(-1,\infty)\setminus\{0\}\) and \(C\) the
 normalizing constant. Then
 \(\int_{-a_1}^{a_2}(f'/f)^2f\,dx<\infty\) if and only if \(m_1>1\) and
 \(m_2>1\).
\end{enumerate}
\end{hbtheorem}

Item~2 says that the location information of a Type~I curve is infinite as
soon as either end exponent is at most \(1\), which is the upper boundary of
fig.~3; item~1 says that for an end exponent below \(1\) the extreme
observation approaches its end faster than \(n^{-1/2}\), so its error
becomes negligible beside that of the mean, as Fisher argues
\cite[p.~348]{Fisher1922Foundations}. For the Type~III curve, item~1 is
item~1 of \cref{thm:p15-w02b-p15-057-endpoint}. At an end exponent equal
to \(1\) the extreme observation still has error of order \(n^{-1/2}\), so
this theorem does not establish that the mean is then of zero efficiency;
\cref{thm:p15-w02b-p15-057-linear-end} does, for a Type~I curve of known
form with one end exponent equal to \(1\), by a different statistic.

\begin{hbtheorem}{thm:p15-w02b-p15-057-rectangle}{proved=annex:p15-w02b-p15-057.rectangle}{The rectangular law: mean against the mean of the extremes}
Let the observations be independent and uniform on
\((m-\frac a2,m+\frac a2)\), \(a>0\), \(n\geq2\). Let \(\overline x\) be
their mean and \(\hat m\) the mean of the least and greatest.
\begin{enumerate}
\item \(E\overline x=E\hat m=m\), \(\operatorname{Var}\overline x
 =a^2/(12n)\), and \(\hat m-m\) has the density
 \(\frac na\bigl(1-\frac{2|x|}a\bigr)^{n-1}\) on \(|x|<\frac a2\), so that
 \(\operatorname{Var}\hat m=a^2/\bigl(2(n+1)(n+2)\bigr)\).
\item \(\operatorname{Var}\hat m/\operatorname{Var}\overline x
 =6n/\bigl((n+1)(n+2)\bigr)\), which is \(6/n+O(n^{-2})\).
\item \(P\bigl(\frac{2n}a(\hat m-m)\leq t\bigr)\) tends to the double
 exponential distribution function with density \(\frac12e^{-|t|}\).
 Treated as a location curve, the double exponential law with density
 \(\frac na e^{-2n|x|/a}\) has intrinsic accuracy \(4n^2/a^2\), and the
 normal law of variance \(a^2/(12n)\) has \(12n/a^2\); their ratio is
 \(3/n\).
\end{enumerate}
\end{hbtheorem}

Item~1 replaces Fisher's large-sample formulas with exact ones: he treats
the two end differences as independent and obtains the double exponential
law and the variance \(a^2/(2n^2)\), hence the ratio \(6/n\)
\cite[pp.~348--349]{Fisher1922Foundations}. His experiment takes
twenty-five samples of \(100\) three-digit numbers from Vega's logarithm
tables, with \(a=1000\) \cite[p.~350]{Fisher1922Foundations}. The exact
standard deviation of \(\hat m\) is then
\(1000/\sqrt{2\cdot101\cdot102}=6.967\) to three decimals, against
Fisher's \(\sqrt{50}=7.07\). Recomputed from his twenty-five printed errors,
the root mean square about the true mean is \(\sqrt{1208.75/25}=6.953\) to
three decimals; his stated \(6.86\) is the root mean square about the
average of the twenty-five errors (\(-1.14\)), which estimates a spread
about that average rather than the error about \(m\). Fisher calls the
6 per cent efficiency conventional, because the rectangle lies outside
the region of validity and the two sampling curves have different forms
\cite[p.~350]{Fisher1922Foundations}; the ratio \(3/n\) of intrinsic
accuracies is his proposed substitute \cite[p.~351]{Fisher1922Foundations}.
Neither ratio is an efficiency in the sense of
\cref{def:p15-w02b-p15-057-efficiency}, since \(\hat m\) has no normal
limit at the \(\sqrt n\) scale.

For the form of the curve Fisher sets up the four-parameter Hessian of
\(-L\) for Type~IV in \((m,a,\nu,r)\), with the transcendental function
\(\mathrm F=e^{-\frac12\nu\pi}\int_0^\pi e^{\nu\theta}\sin^r\theta\,d\theta\),
specializes it to Type~VII (\(\nu=0\)), where it splits into two
\(2\times2\) factors, and compares the resulting limit variances of the
optimum estimates of \(\nu\) and \(r\) with those of the moment estimates
from \(\beta_1\) and \(\beta_2\) \cite[pp.~351--353]{Fisher1922Foundations};
the same is done for Type~II \cite[pp.~354--355]{Fisher1922Foundations}.
In those pages his function written \(F(x)\) with an argument such as
\(r/2\) is the function \(F\) of \cref{sec:p15-w02b-p15-057-type-iii},
\(F(x)=\psi'(x+1)\); it is written here through \(\psi'\) to keep it apart
from the normalizing integral \(\mathrm F\). Fisher does not define the
moment estimates beyond naming them; the definition used below solves the
two moment equations locally, which is the only reading under which they
are functions of \(\sqrt{b_1}\) and \(b_2\) near a symmetric curve.

\begin{hbdefinition}{def:p15-w02b-p15-057-form-moments}{Moment estimates of the form near a symmetric curve}
Let \(m_2,m_3,m_4\) be the sample central moments of \(x_1,\ldots,x_n\), and
put \(g_1=m_3/m_2^{3/2}\) and \(b_2=m_4/m_2^2\) when \(m_2>0\). Let a family
of curves be indexed by form parameters \((\nu,r)\) on an open set
\(U\subseteq\mathbb R^2\) together with location and scale, and let
\(\gamma_1(\nu,r)\), \(\beta_2(\nu,r)\) be the skewness
\(\mu_3/\mu_2^{3/2}\) and kurtosis \(\mu_4/\mu_2^2\) of one observation.
Suppose that at \((0,r_0)\in U\) the map
\(\Gamma=(\gamma_1,\beta_2)\) is continuously differentiable with
nonsingular derivative, and let \(\Psi=(\Psi_1,\Psi_2)\) be its
continuously differentiable inverse from an open neighbourhood \(W\) of
\(\Gamma(0,r_0)\) onto an open neighbourhood of \((0,r_0)\). The
\emph{moment estimates} are \((\nu_\mu,r_\mu)=\Psi(g_1,b_2)\) when
\(m_2>0\) and \((g_1,b_2)\in W\), and \((0,r_0)\) otherwise.
\end{hbdefinition}

\begin{hbtheorem}{thm:p15-w02b-p15-057-form-vii}{proved=annex:p15-w02b-p15-057.form-vii}{The form parameters of the Type IV and Type VII curves}
For \(m\in\mathbb R\), \(a>0\), \(\nu\in\mathbb R\) and \(r>-1\), let
\(\xi=(x-m)/a\) and
\[
 f(x)=\frac{1}{a\,\mathrm F(\nu,r)}\,(1+\xi^2)^{-(r+2)/2}e^{-\nu\tan^{-1}\xi},
 \qquad
 \mathrm F(\nu,r)=e^{-\frac12\nu\pi}\int_0^\pi e^{\nu\theta}\sin^r\theta\,d\theta,
\]
and let the observations be independent with density \(f\). Put
\(D=(r+4)^2+\nu^2\) and \(D_2=(r+2)^2+\nu^2\).
\begin{enumerate}
\item \(\mathrm F\) is the normalizing constant, and the model with
 parameter \((m,a,\nu,r)\in\mathbb R\times(0,\infty)\times\mathbb R
 \times(-1,\infty)\) is regular in the sense of
 \cref{def:p15-w02b-p15-057-multi-regular} at every point. Its
 per-observation information matrix is
 \[
  J=\begin{pmatrix}
  \frac{(r+1)(r+2)(r+4)}{a^2D}&-\frac{\nu(r+1)(r+2)}{a^2D}&
   -\frac{(r+1)(r+2)}{aD_2}&\frac{\nu(r+1)}{aD_2}\\[3pt]
  -\frac{\nu(r+1)(r+2)}{a^2D}&\frac{(r+1)(2(r+4)+\nu^2)}{a^2D}&
   \frac{\nu(r+1)}{aD_2}&-\frac{r+2+\nu^2}{aD_2}\\[3pt]
  -\frac{(r+1)(r+2)}{aD_2}&\frac{\nu(r+1)}{aD_2}&
   \partial_\nu^2\log\mathrm F&\partial_\nu\partial_r\log\mathrm F\\[3pt]
  \frac{\nu(r+1)}{aD_2}&-\frac{r+2+\nu^2}{aD_2}&
   \partial_\nu\partial_r\log\mathrm F&\partial_r^2\log\mathrm F
  \end{pmatrix},
 \]
 and it is positive definite.
\item At \(\nu=0\), \(\partial_\nu\partial_r\log\mathrm F=0\),
 \(\partial_\nu^2\log\mathrm F=\tfrac12\psi'(\tfrac r2+1)\) and
 \(\partial_r^2\log\mathrm F=\tfrac14\bigl(\psi'(\tfrac{r+1}2)
 -\psi'(\tfrac r2+1)\bigr)\), so \(J\) is block diagonal in the pairs
 \((m,\nu)\) and \((a,r)\). Put
 \[
  G_\nu(r)=(r+2)^3\psi'(\tfrac r2+1)-2(r+1)(r+4),\quad
  G_r(r)=(r+1)^2(r+2)^2\bigl(\psi'(\tfrac{r+1}2)-\psi'(\tfrac r2+1)\bigr)
  -2(r+1)(r+4).
 \]
 Both are positive, and the \((\nu,\nu)\) and \((r,r)\) entries of
 \(J^{-1}\) are \(2(r+2)^3/G_\nu(r)\) and \(4(r+1)^2(r+2)^2/G_r(r)\). By
 \cref{thm:p15-w02b-p15-057-multi-regular}, every sequence of roots of the
 likelihood equations converging in probability to \((m,a,0,r)\) has
 \(\sqrt n\,\hat\nu\) and \(\sqrt n(\hat r-r)\) asymptotically normal with
 these variances.
\item Let \(r_0>3\) and \(\nu=0\). In the Type~IV family,
 \(\gamma_1=-4\nu\sqrt{r-1}\big/\bigl((r-2)\sqrt{r^2+\nu^2}\bigr)\) and
 \(\beta_2\) of item~4 of \cref{thm:p15-w02b-p15-057-typeiv} satisfy the
 hypotheses of \cref{def:p15-w02b-p15-057-form-moments} at \((0,r_0)\).
 If \(r_0>5\),
 \[
  \sqrt n\,\nu_\mu\Rightarrow N\Bigl(0,\ \frac38\cdot
  \frac{r_0^2(r_0-2)^2(r_0^2+r_0+10)}{(r_0-1)(r_0-3)(r_0-5)}\Bigr),
 \]
 and if \(r_0>7\), \(\sqrt n\,(r_\mu-r_0)\) converges in law to
 \[
  N\Bigl(0,\ \frac23\cdot
  \frac{r_0(r_0-1)^2(r_0-3)(r_0^2-r_0+18)}{(r_0-5)(r_0-7)}\Bigr).
 \]
\item The efficiencies of \(\nu_\mu\) and \(r_\mu\) in the sense of
 \cref{def:p15-w02b-p15-057-multi-efficiency} at \((m,a,0,r)\) are
 \[
  E_\nu(r)=\frac{16(r+2)^3(r-1)(r-3)(r-5)}
  {3\,G_\nu(r)\,r^2(r-2)^2(r^2+r+10)}\quad(r>5),\qquad
  E_r(r)=\frac{6(r+1)^2(r+2)^2(r-5)(r-7)}
  {G_r(r)\,r(r-1)^2(r-3)(r^2-r+18)}\quad(r>7).
 \]
 Moreover \(G_\nu(r)=\frac{16}3-\frac{16}{15(r+2)^2}
 +\frac{64\vartheta}{21(r+2)^4}\) for some \(\vartheta\in(0,1)\),
 \(G_r(r)=6-\frac4{(r+2)^2}+O(r^{-3})\), and as \(r\to\infty\)
 \[
  E_\nu(r)=1-\frac{144}{5r^2}+O(r^{-3}),\qquad
  E_r(r)=1-\frac{160}{3r^2}+O(r^{-3}).
 \]
\end{enumerate}
\end{hbtheorem}

Item~1 is Fisher's determinant \cite[p.~351]{Fisher1922Foundations} with
one difference of sign: his four entries pairing \(m\) or \(a\) with
\(\nu\) or \(r\) have the opposite signs. His matrix is \(\Lambda J\Lambda\)
with \(\Lambda=\operatorname{diag}(1,1,-1,-1)\), which has the same
determinant and the same diagonal entries of the inverse, so his standard
deviations are unaffected; the correlations of \(\hat m\) and \(\hat a\)
with \(\hat\nu\) and \(\hat r\) read from his matrix have the wrong sign.
Item~2 is his Type~VII determinant, its two factors, and his
\(n\sigma_\nu^2=2(r+2)^3/G_\nu\) and
\(n\sigma_r^2=4(r+1)(r+2)^2/\bigl(G_r/(r+1)\bigr)\)
\cite[p.~352]{Fisher1922Foundations}. Item~3 gives his
\(n\sigma_{\nu_\mu}^2\) and \(n\sigma_{r_\mu}^2\)
\cite[p.~352]{Fisher1922Foundations}, now with the moment conditions they
need: the first requires a finite sixth moment (\(r>5\)) and the second a
finite eighth moment (\(r>7\)), by item~1 of
\cref{thm:p15-w02b-p15-057-pearson-moments}. Item~4 contains his
approximations \(G_\nu\approx\frac{16}3(1-\frac1{5(r+2)^2})\) and
\(G_r\approx6-\frac4{(r+2)^2}\) and his large-\(r\) efficiencies
\(1-28.8/r^2\) and \(1-53.3/r^2\) \cite[pp.~352--353]{Fisher1922Foundations},
with \(144/5=28.8\) and \(160/3=53.\dot3\); the first approximation is
exact up to the stated term in \(\vartheta\).

Evaluating item~4 with the trigamma function reproduces Fisher's Type~VII
table \cite[p.~353]{Fisher1922Foundations} for \(r=5,\ldots,18\) to within
three units in the last printed place, with one exception: at \(r=10\) the
efficiency of \(r_\mu\) is \(0.4289\), not the printed \(0.4403\). His
entries \(0\) at \(r=5\) for \(\nu_\mu\) and at \(r=7\) for \(r_\mu\) are
limits: \(E_\nu(r)\to0\) as \(r\downarrow5\) and \(E_r(r)\to0\) as
\(r\downarrow7\), and at and below those values the moment estimate has no
normal limit at the \(\sqrt n\) scale, so no efficiency is defined. The
efficiency of \(r_\mu\) equals \(0.8\) at \(r=17.172\), where
\(\beta_2=3(r-1)/(r-3)=3.4234\), and that of \(\nu_\mu\) equals \(0.8\) at
\(r=12.112\) (three decimals). Since \(E_\nu>E_r\) at each tabulated
\(r\), the first of these is Fisher's statement that 80 per cent is not
reached for the form until \(r\) is about \(17.2\), with \(\beta_2=3.42\)
\cite[p.~353]{Fisher1922Foundations}.

For Type~II, Fisher writes \(r\) for the positive quantity \(-r\)
\cite[p.~354]{Fisher1922Foundations}. The curve is then of finite range,
and the asymmetric curves next to it are of Type~I; the next theorem uses
the parametrization of those curves for which Fisher's Type~II matrix is
the information matrix.

\begin{hbtheorem}{thm:p15-w02b-p15-057-form-ii}{proved=annex:p15-w02b-p15-057.form-ii}{The form parameters of the Type II curves}
For \(m\in\mathbb R\), \(a>0\), \(r>0\) and \(|\nu|<r\), let
\(\xi=(x-m)/a\) and
\[
 f(x)=\frac1{aZ(\nu,r)}(1+\xi)^{(r-2-\nu)/2}(1-\xi)^{(r-2+\nu)/2}
 \quad(|\xi|<1),\qquad
 Z(\nu,r)=\frac{2^{r-1}\Gamma(\frac{r-\nu}2)\Gamma(\frac{r+\nu}2)}{\Gamma(r)},
\]
and \(f=0\) otherwise; \(\nu=0\) is the Type~II curve
\(\propto(1-\xi^2)^{(r-2)/2}\). Let the observations be independent with
density \(f\) at \((m,a,0,r)\).
\begin{enumerate}
\item If \(r>4\), condition (S) of
 \cref{def:p15-w02b-p15-057-multi-conditions} holds at \((m,a,0,r)\), and
 the information matrix is block diagonal in \((m,\nu)\) and \((a,r)\)
 with
 \[
  \begin{pmatrix}\frac{(r-1)(r-2)}{a^2(r-4)}&-\frac{r-1}{a(r-2)}\\
  -\frac{r-1}{a(r-2)}&\frac12\psi'(\frac r2)\end{pmatrix},\qquad
  \begin{pmatrix}\frac{2(r-1)}{a^2(r-4)}&-\frac1{a(r-2)}\\
  -\frac1{a(r-2)}&\frac14\bigl(\psi'(\frac r2)-\psi'(\frac{r+1}2)\bigr)
  \end{pmatrix}.
 \]
 With \(H_\nu(r)=(r-2)^3\psi'(\frac r2)-2(r-1)(r-4)\) and
 \(H_r(r)=(r-1)^2(r-2)^2\bigl(\psi'(\frac r2)-\psi'(\frac{r+1}2)\bigr)
 -2(r-1)(r-4)\), both positive, the inverse-information values for
 \(\nu\) and \(r\) are \(2(r-2)^3/H_\nu(r)\) and
 \(4(r-1)^2(r-2)^2/H_r(r)\). Moreover \(H_\nu(r)=G_\nu(r-4)\) and
 \(H_r(r)=G_r(r-3)\), with \(G_\nu,G_r\) as in
 \cref{thm:p15-w02b-p15-057-form-vii}.
\item For every \(r_0>0\), the skewness and kurtosis of this family
 satisfy the hypotheses of \cref{def:p15-w02b-p15-057-form-moments} at
 \((0,r_0)\), with
 \(\gamma_1=4\nu\sqrt{r+1}\big/\bigl((r+2)\sqrt{r^2-\nu^2}\bigr)\) and
 \(\beta_2(0,r)=3(r+1)/(r+3)\), and
 \[
  \sqrt n\,\nu_\mu\Rightarrow N\Bigl(0,\frac38\cdot
  \frac{r_0^2(r_0+2)^2(r_0^2-r_0+10)}{(r_0+1)(r_0+3)(r_0+5)}\Bigr),
 \]
 and \(\sqrt n\,(r_\mu-r_0)\) converges in law to
 \[
  N\Bigl(0,\frac23\cdot
  \frac{r_0(r_0+1)^2(r_0+3)(r_0^2+r_0+18)}{(r_0+5)(r_0+7)}\Bigr).
 \]
\item For \(r>4\) the efficiencies of \(\nu_\mu\) and \(r_\mu\) in the sense
 of \cref{def:p15-w02b-p15-057-multi-efficiency} are the ratios of the
 values in item~1 to those in item~2, and as \(r\to\infty\) they are
 \(1-\frac{144}{5r^2}+O(r^{-3})\) and \(1-\frac{160}{3r^2}+O(r^{-3})\).
\end{enumerate}
\end{hbtheorem}

Item~1 is Fisher's Type~II Hessian and his \(n\sigma_\nu^2\),
\(n\sigma_r^2\) \cite[p.~354]{Fisher1922Foundations}; its signs agree with
his. The identity \(H_\nu(r)=G_\nu(r-4)\) is his remark that the Type~II
expression is the same function of \(r-4\) as the Type~VII one is of
\(r\), obtained from \(\psi'(\frac r2)=\psi'(\frac r2-1)-4/(r-2)^2\)
\cite[p.~354]{Fisher1922Foundations}; \(H_r(r)=G_r(r-3)\) is his second
such identity, whose left side is printed with the factor
\((r+1)^2\) where his table heading and the identity itself require
\((r-1)^2\). Fisher prints no Type~II moment variances; item~2 supplies
the ones behind his table, and they are the expressions of item~3 of
\cref{thm:p15-w02b-p15-057-form-vii} with \(r\) replaced by \(-r\), in
line with his remark that \(r\) changes sign between the two types
\cite[p.~355]{Fisher1922Foundations}. Every moment is finite on a finite
range, so there is no lower moment threshold. Evaluating item~3
reproduces his Type~II table for \(r=5,\ldots,15\) to within two units
in the last printed place \cite[p.~355]{Fisher1922Foundations}. His
entries at \(r=2,3,4\) (\(0\), \(0.0576\), \(0.2056\) for \(\nu_\mu\) and
\(0\), \(0.0431\), \(0.1445\) for \(r_\mu\)) are values of the same
formulas where the location entry of the information is infinite
(\(r\leq4\), the end exponent \((r-2)/2\) being at most \(1\), as in
item~2 of \cref{thm:p15-w02b-p15-057-terminus}); they are not information
ratios. The efficiency of \(r_\mu\) equals \(0.8\) at \(r=14.059\), where
\(\beta_2=3(r+1)/(r+3)=2.6483\), which is Fisher's \(r\) about \(14.1\)
with \(\beta_2=2.65\) \cite[p.~355]{Fisher1922Foundations}; that of
\(\nu_\mu\) equals \(0.8\) at \(r=10.985\). His conclusion that for
symmetric Pearson curves the moments have at least 80 per cent efficiency
for the form exactly when \(\beta_2\) lies between \(2.65\) and \(3.42\)
\cite[p.~355]{Fisher1922Foundations} is the statement that \(E_r\) is
increasing in \(r\) on each side, which the tables support at the
tabulated points but which is not proved here. The Type~II efficiencies
are inverse-information ratios under (S); because the range
\((m-a,m+a)\) moves with the parameter,
\cref{thm:p15-w02b-p15-057-multi-regular} does not apply, and condition
(C), which would make the inverse-information values the limit variances
of a root of the likelihood equations, is not verified here.

\section{Why moments are efficient near the normal curve}
\label{sec:p15-w02b-p15-057-near-normal}

Fisher's Section~11 explains the high efficiency of moments near the normal
point by the family of curves for which the method of moments is the
method of maximum likelihood, namely densities whose logarithm is a
polynomial of degree four \cite[pp.~355--356]{Fisher1922Foundations}.

\begin{hbtheorem}{thm:p15-w02b-p15-057-quartic}{proved=annex:p15-w02b-p15-057.quartic}{The quartic-exponent family}
For \(\theta=(\theta_1,\ldots,\theta_4)\in\mathbb R^4\) put
\(T(x)=(x,x^2,x^3,x^4)\) and \(Z(\theta)=\int_{\mathbb R}e^{\theta\cdot
T(x)}\,dx\in(0,\infty]\).
\begin{enumerate}
\item \(Z(\theta)<\infty\) if and only if \(\theta_4<0\), or
 \(\theta_4=\theta_3=0\) and \(\theta_2<0\).
\item On the open set \(\Theta^\circ=\{\theta_4<0\}\), let
 \(f_\theta=e^{\theta\cdot T-\log Z(\theta)}\). Then \(\log Z\) is
 differentiable with gradient \(E_\theta T(X)\), and for any observations
 \(x_1,\ldots,x_n\) the log likelihood
 is strictly concave on \(\Theta^\circ\); its gradient vanishes at
 \(\theta\) if and only if \(E_\theta X^j=n^{-1}\sum_ix_i^j\) for
 \(j=1,2,3,4\). Such a \(\theta\) is unique and is the maximum likelihood
 estimate on \(\Theta^\circ\).
\item Let \(\sigma>0\), \(k_1,k_2\in\mathbb R\) and
 \(\phi(x)=-\frac{x^2}{2\sigma^2}+k_1\frac{x^3}{\sigma^3}
 +k_2\frac{x^4}{\sigma^4}\). Then \(e^\phi\) is integrable if and only if
 \(k_2<0\) or \(k_1=k_2=0\). With
 \(u=3k_1x/\sigma+4k_2x^2/\sigma^2\) and \(u\neq-1\),
 \[
  \phi'(x)=-\frac{x}{\sigma^2(1+u)}+\frac{x\,u^2}{\sigma^2(1+u)} ,
 \]
 so on \(|x|\leq L\sigma\) and for \(|k_1|+|k_2|\leq1/(2(3L+4L^2))\),
 \(\bigl|\phi'(x)+x/(\sigma^2(1+u))\bigr|\leq
 2L(3L+4L^2)^2(|k_1|+|k_2|)^2/\sigma\).
\end{enumerate}
\end{hbtheorem}

Item~2 is the exact form of Fisher's deduction that \(\partial_\theta\log
y\) must be a polynomial of degree at most four in \(x\), with the
normalization tying the five coefficients
\cite[p.~356]{Fisher1922Foundations}: the first four sample moments are the
maximum likelihood estimates in that family, and, with \(\theta_3=\theta_4=0\),
the first two are the maximum likelihood estimates for the normal curve.
Item~3 corrects and bounds the comparison that follows. Fisher writes the
near-normal quartic curve \(Ce^{\phi}\) without a sign condition; it is a
frequency curve only for \(k_2<0\) (or the normal curve itself). Its
logarithmic derivative agrees with Pearson's form
\(-x/(\sigma^2(1+u))\), whose denominator is quadratic in \(x\), up to a
term of second order in \((k_1,k_2)\), but only uniformly on a bounded
range of \(x/\sigma\). In the tails the two differ: for \(k_2<0\) the
quadratic \(1+u\) has two real zeros, one on each side of the origin, so
for small \(|k_1|+|k_2|\) the Pearson curve with
logarithmic derivative \(-x/(\sigma^2(1+u))\) has a finite range ending at
such a zero, whereas the quartic curve is positive on the whole line. The
agreement is therefore local, and it supports Fisher's
attribution of the efficiency of \(\mu_3\) and \(\mu_4\) near the normal
point \cite[p.~356]{Fisher1922Foundations} as an explanation, not as a proof
that those efficiencies tend to one.

\section{The end point of the Type III curve}
\label{sec:p15-w02b-p15-057-endpoint}

For \(p<1\) Fisher states that the location of the Type III curve is
determined by its end point, and that from \(p=1\) downwards the weight of
that determination grows faster than the sample size
\cite[p.~337]{Fisher1922Foundations}. The following theorem gives the
exact rate at which the least observation approaches the end point, with
or without knowledge of \(a\) and \(p\).

\begin{hbtheorem}{thm:p15-w02b-p15-057-endpoint}{proved=annex:p15-w02b-p15-057.endpoint}{The least observation of a Type III sample}
Let the observations be independent with the Type III law of
\cref{sec:p15-w02b-p15-057-type-iii}, \(a>0\), \(p>-1\), and let
\(X_{(1)}=\min_{i\leq n}X_i\).
\begin{enumerate}
\item For every \(t\geq0\),
 \[
  P\Bigl(n^{1/(p+1)}\,\frac{X_{(1)}-m}{a}>t\Bigr)\to
  \exp\Bigl(-\frac{t^{p+1}}{\Gamma(p+2)}\Bigr).
 \]
\item If \(-1<p<1\), then \(\sqrt n\,(X_{(1)}-m)\to0\) in probability. If
 \(T_n\) is any estimator sequence with
 \(\sqrt n\,(T_n-m)\) converging in law to \(N(0,v)\) and \(v>0\), in particular the
 moment estimate \(m_\mu\) of \cref{thm:p15-w02b-p15-057-typeiii} with
 \(v=a^2(p+1)\), then
 \(P\bigl(|X_{(1)}-m|<|T_n-m|\bigr)\to1\).
\item If \(p=1\), then \(\sqrt n\,(X_{(1)}-m)/a\) converges in law to the
 distribution with \(P(\,\cdot>t)=e^{-t^2/2}\) for \(t\geq0\). If \(p>1\),
 then \(\sqrt n\,(X_{(1)}-m)\to\infty\) in probability.
\end{enumerate}
\end{hbtheorem}

Item~2 is the exact content of the end-point regime for \(-1<p<1\): no
\(\sqrt n\)-normal location estimate can compete with the least
observation, so a comparison of variances at the \(\sqrt n\) scale, and
with it the efficiency of \cref{def:p15-w02b-p15-057-efficiency}, has no
application there. Item~3 shows that the least observation alone does not
supply the faster-than-\(n\) weight Fisher asserts at \(p=1\) itself: its
error is of exact order \(n^{-1/2}\) there, with a non-normal limit. Another
statistic does achieve a faster rate at \(p=1\) when \(a\) is known; the
next theorem constructs it for every location family whose density
vanishes linearly at a finite end, which is what Fisher's ``finite angle
with the axis'' means \cite[p.~347]{Fisher1922Foundations}.

\begin{hbtheorem}{thm:p15-w02b-p15-057-linear-end}{proved=annex:p15-w02b-p15-057.linear-end}{Location faster than \(\sqrt n\) at an end exponent equal to one}
Let \(f\) be a known probability density on \(\mathbb R\) with \(f(y)=0\)
for \(y\leq0\), and suppose there are \(0<c\leq C<\infty\) and
\(\delta>0\) with \(cy\leq f(y)\leq Cy\) for \(0<y<\delta\). Let
\(X_i=m+Y_i\), where \(Y_1,Y_2,\ldots\) are independent with density
\(f\) and \(m\in\mathbb R\) is unknown. For \(0<\varepsilon<\delta\) put
\[
 g_\varepsilon(y)=\begin{cases}1/\varepsilon,&y\leq\varepsilon,\\
 1/y,&y>\varepsilon,\end{cases}\qquad
 \kappa_\varepsilon=E\,g_\varepsilon(Y_1),\qquad
 J(\varepsilon)=\int_\varepsilon^\infty\frac{f(y)}{y^2}\,dy,
\]
\(\psi_\varepsilon=\kappa_\varepsilon-g_\varepsilon\) and
\(\rho_\varepsilon(y)=\int_0^y\psi_\varepsilon(u)\,du\).
\begin{enumerate}
\item \(\rho_\varepsilon\) is a criterion of
 \cref{def:p15-convex-location-criterion} with continuous
 \(\psi=\psi_\varepsilon\). The midpoint \(T_{n,\varepsilon}\) of its
 minimizing set is translation equivariant, and under every \(m\)
 \[
  \sqrt n\,(T_{n,\varepsilon}-m)\Rightarrow N\bigl(0,v(\varepsilon)\bigr),
  \qquad
  v(\varepsilon)=\frac{\operatorname{Var}g_\varepsilon(Y_1)}{J(\varepsilon)^2}
  \leq\frac{C}{2J(\varepsilon)^2}+\frac1{J(\varepsilon)},
 \]
 with \(J(\varepsilon)\geq c\log(\delta/\varepsilon)\); hence
 \(v(\varepsilon)\to0\) as \(\varepsilon\to0\).
\item There are \(\varepsilon_k\downarrow0\) and integers
 \(N_1<N_2<\cdots\), depending on \(f\) alone, such that the
 translation-equivariant estimator \(T_n=T_{n,\varepsilon_k}\) for
 \(N_k\leq n<N_{k+1}\) (and \(T_n=T_{n,\varepsilon_1}\) for \(n<N_1\))
 satisfies \(\sqrt n\,(T_n-m)\to0\) in probability under every \(m\).
\item If \(S_n\) is any estimator sequence with
 \(\sqrt n\,(S_n-m)\Rightarrow N(0,v)\), \(v>0\), then
 \(P_m(|T_n-m|<|S_n-m|)\to1\), and the relative efficiency of \(S\) with
 respect to \(T_{\cdot,\varepsilon}\), \(v(\varepsilon)/v\), tends to
 \(0\) as \(\varepsilon\to0\). If \(\int y^2f(y)\,dy<\infty\), this applies
 to \(S_n=\overline X-\int yf(y)\,dy\), with \(v=\operatorname{Var}Y_1\).
\item The hypothesis holds for the Type~III density with \(p=1\) and known
 \(a\), \(f(y)=a^{-2}ye^{-y/a}\) (with \(c=e^{-1}a^{-2}\), \(C=a^{-2}\),
 \(\delta=a\)), and, with \(m=-a_1\), for the Type~I density of
 \cref{thm:p15-w02b-p15-057-terminus} with known \(a_1,a_2,m_2\) and
 \(m_1=1\); for \(m_2=1\) apply it to \(-X_i\).
\end{enumerate}
\end{hbtheorem}

The estimator downweights observations far from the end as \(1/y\) and
caps the weight near it; the cap is what makes the criterion convex and the
limit variance finite, and lowering the cap makes that variance as small as
desired, because the integral \(J(\varepsilon)\) diverges logarithmically.
Thus at \(p=1\) the weight of the location determination does grow faster
than \(n\), as Fisher states \cite[p.~337]{Fisher1922Foundations}, in the
exact sense \(n\,(T_n-m)^2\to0\) in probability, and the mean has
efficiency zero relative to the attainable \(\sqrt n\)-normal estimates, as
he states for the curves meeting the axis at a finite angle
\cite[pp.~347--348]{Fisher1922Foundations}. Three limits remain. The
theorem assumes \(a\) (and the shape) known, whereas Fisher's sentence on
p.~337 is meant whether or not \(p\) and \(a\) are known; the unknown-scale
case is not proved here. The rate obtained is only
\(o(n^{-1/2})\): the bound \(v(\varepsilon)\approx1/\log(1/\varepsilon)\)
suggests error of order \((n\log n)^{-1/2}\), and that order is not proved
here. And the diagonal estimator is a device of proof, whose blocks
\(N_k\) are not computed.

\begin{hbexample}{def:p15-w02b-p15-057-linear-end-example}{The Type III curve at \(p=1\) with \(a=1\)}
Here \(f(y)=ye^{-y}\), \(\operatorname{Var}Y_1=2\), and every quantity in
item~1 is explicit:
\(J(\varepsilon)=\int_\varepsilon^\infty e^{-y}y^{-1}dy\),
\(\kappa_\varepsilon=\varepsilon^{-1}\bigl(1-e^{-\varepsilon}(1+\varepsilon)\bigr)+e^{-\varepsilon}\)
and
\(E\,g_\varepsilon(Y_1)^2=\varepsilon^{-2}\bigl(1-e^{-\varepsilon}(1+\varepsilon)\bigr)+J(\varepsilon)\).
At \(\varepsilon=0.01\), \(J=4.03793\), \(\kappa=0.99502\) and
\(v(0.01)=0.21739\), so the mean has relative efficiency
\(0.21739/2=0.1087\) with respect to \(T_{\cdot,0.01}\); at
\(\varepsilon=10^{-6}\), \(J=13.23830\) and \(v=0.07269\), relative
efficiency \(0.0363\). The decrease is slow, as the logarithm predicts.
\end{hbexample}

\section{Grouped normal observations}
\label{sec:p15-w02b-p15-057-grouped}

\Cref{sec:p15-fisher-grouping} sets up the class likelihood for normal
observations and leaves three questions open: whether the grouped
maximum likelihood estimate exists and what its sampling limit is, how
far the information at one fixed grid phase can be from its phase average,
and how the discrepancy between observed and fitted counts is calibrated.
The first two are answered below under stated hypotheses.

\begin{hbdefinition}{def:p15-w02b-p15-057-grouped}{The finite grouped normal experiment}
Fix \(K\geq3\) and endpoints \(b_1<\cdots<b_{K-1}\), and put
\(b_0=-\infty\), \(b_K=\infty\). For \(\theta=(\mu,\sigma)\in
\Theta=\mathbb R\times(0,\infty)\) the cell probabilities are
\(p_j(\theta)=\Phi(z_{j+1})-\Phi(z_j)\), \(z_j=(b_j-\mu)/\sigma\),
\(j=0,\ldots,K-1\), as in \cref{def:p15-fg-normal}. Each observation
records only the cell into which an independent normal variable with
parameters \(\theta_0\) falls; \(n_j\) is the count of cell \(j\) among
the first \(n\) observations, and the grouped log likelihood is
\(\ell_n(\theta)=\sum_jn_j\log p_j(\theta)\). Cells \(0\) and \(K-1\) are
the tail cells; the others are bounded. Two cells \(j<j'\) are adjacent
when \(j'=j+1\).
\end{hbdefinition}

\begin{hbtheorem}{thm:p15-w02b-p15-057-grouped-mle}{proved=annex:p15-w02b-p15-057.grouped-mle}{Existence, consistency and limit of the grouped estimate}
In the experiment of \cref{def:p15-w02b-p15-057-grouped}:
\begin{enumerate}
\item If the occupied cells include a bounded cell and two cells that are
 not adjacent, \(\ell_n\) attains its supremum on \(\Theta\). If all
 observations lie in one cell, or in two adjacent cells, it does not.
 The first condition holds, almost surely, for all large \(n\).
\item \(\theta\) is identified by the cell probabilities, and outside one
 null event every choice of maximizers \(\widehat\theta_n\) satisfies
 \(\widehat\theta_n\to\theta_0\).
\item The matrix
 \(J=\sum_jp_j(\theta_0)^{-1}\nabla p_j(\theta_0)\nabla
 p_j(\theta_0)^{\top}\) is positive definite, and every sequence of random
 vectors that are maximizers whenever a maximizer exists satisfies
 \(\sqrt n\,t^{\top}(\widehat\theta_n-\theta_0)\Rightarrow
 N(0,t^{\top}J^{-1}t)\) for each \(t\in\mathbb R^2\), with
 \(-n^{-1}\nabla^2\ell_n(\widehat\theta_n)\to J\) in probability.
\end{enumerate}
\end{hbtheorem}

The theorem concerns a finite partition with both tails kept as cells, as
\cref{sec:p15-fisher-grouping} requires for the Smith data, and normal
sampling; it does not cover a model fitted to counts from some other law.
Item~1 is a sample-by-sample statement: with positive probability at each
\(n\) all observations fall into two adjacent cells, and then the
likelihood approaches its supremum as \(\sigma\to0\) without attaining
it. The matrix \(J\) is the grouped information of
\cref{thm:p15-fg-information}, so the limit variance of each linear
combination is at least its ungrouped value.

\begin{hbtheorem}{thm:p15-w02b-p15-057-grouped-phase}{proved=annex:p15-w02b-p15-057.grouped-phase}{A fixed-phase bound for grouped normal information}
Let \(X\) be normal with parameters \(\mu,\sigma\), grouped on the grid of
\cref{def:p15-fg-grid} with width \(h\) and phase \(\vartheta\), and put
\(a=h/\sigma\). Let \(I_\mu(\vartheta)\) and \(I_\sigma(\vartheta)\) be the
location and scale information of the recorded class, each with the other
parameter known, and \(\overline I_\mu,\overline I_\sigma\) their averages
over a uniform phase. Let \(Z\) be standard normal. Then for every phase
and every \(a\in(0,2]\),
\[
 \Bigl|\sigma^2I_\mu(\vartheta)-1+\frac{a^2}{12}\Bigr|\leq
 \varepsilon_\mu(a)=\frac{a^2}{12}\bigl(2e^{3a^2/2}\Phi(a)-1\bigr),
\]
\[
 \Bigl|\sigma^2I_\sigma(\vartheta)-2+\frac{a^2}{3}\Bigr|\leq
 \varepsilon_\sigma(a)=\frac{a^2}3\Bigl(a+\frac{a^2}4\Bigr)+\frac{a^4}{180}
 +E\Bigl[\Bigl(\frac{a^2(|Z|+a/2)^2}{3}+\frac{a^4}{180}\Bigr)
 \bigl(e^{a(|Z|+a)}-1\bigr)\Bigr].
\]
Consequently the periodic remainders satisfy
\(\sigma^2|I_\mu(\vartheta)-\overline I_\mu|\leq2\varepsilon_\mu(a)\) and
\(\sigma^2|I_\sigma(\vartheta)-\overline I_\sigma|\leq2\varepsilon_\sigma(a)\),
and \(\varepsilon_\mu(a)/a^3\) and \(\varepsilon_\sigma(a)/a^3\) stay bounded
as \(a\to0\).
\end{hbtheorem}

The bound is uniform in phase and makes no small-width approximation: at
every fixed grid the losses of location and scale information are
\(a^2/12\) and \(a^2/6\) of their ungrouped values up to an explicit error
of order \(a^3\). It does not show that the periodic part is as small as the
exponentially small moment harmonics of \cref{thm:p15-fg-fourier}, and so
by itself it does not decide the eighth-order terms of
\cref{thm:p15-fg-expansion} at a fixed phase. At \(a=1\) the location bound
is about \(0.55\), which says nothing useful. The periodic part is in fact
exponentially small, by the same Poisson summation that gives the moment
harmonics, applied to the nonlinear cell-score integrands after they are
continued into a strip of the complex plane.

\begin{hblemma}{thm:p15-w02b-p15-057-periodization}{proved=annex:p15-w02b-p15-057.periodization}{Periodization and its Fourier series}
Let \(g:\mathbb R\to\mathbb C\) be continuous, let \(L>0\), and suppose
that \(\sum_{j\in\mathbb Z}\sup_{0\leq u\leq L}|g(u+jL)|<\infty\). Put
\(F(u)=\sum_{j\in\mathbb Z}g(u+jL)\),
\(\widehat g(\omega)=\int_{\mathbb R}g(x)e^{-i\omega x}dx\) and
\(\omega_k=2\pi k/L\). Then \(F\) is continuous with period \(L\), its
average over a period is \(L^{-1}\widehat g(0)\), and if
\(\sum_k|\widehat g(\omega_k)|<\infty\), then for every \(u\)
\[
 F(u)=\frac1L\sum_{k\in\mathbb Z}\widehat g(\omega_k)e^{i\omega_ku},
 \qquad
 \Bigl|F(u)-\frac{\widehat g(0)}L\Bigr|\leq\frac1L\sum_{k\neq0}
 |\widehat g(\omega_k)|.
\]
If moreover \(g\) extends to a function analytic on the strip
\(|\operatorname{Im}z|<c'\), and for some \(0<c<c'\) there is \(M\) with
\(\int_{\mathbb R}|g(x+iy)|\,dx\leq M\) for \(|y|\leq c\) and
\(g(x+iy)\to0\) as \(|x|\to\infty\) uniformly in \(|y|\leq c\), then
\(|\widehat g(\omega)|\leq Me^{-c|\omega|}\) for every real \(\omega\).
\end{hblemma}

\begin{hbtheorem}{thm:p15-w02b-p15-057-grouped-exp}{proved=annex:p15-w02b-p15-057.grouped-exp}{Fixed-phase grouped normal information to every order}
In the setting of \cref{thm:p15-w02b-p15-057-grouped-phase}, whose class
endpoints are \(b+jh\), \(j\in\mathbb Z\), with
\(b=h(\vartheta/(2\pi)-1/2)\), let \(I_{\mu\sigma}(\vartheta)\)
be the cross entry of the information matrix of the recorded class, and let
\(J_\mu,J_\sigma\) be the phase-averaged values of
\cref{thm:p15-fg-expansion}. Put \(\beta=\sqrt{2/\pi}\) and, for
\(0<\tau<\pi/2\), \(c=\tau/a\), \(\eta=e^{-2\pi\tau/a^2}\),
\[
 B_\mu=\frac{2e^{c^2/2}(1+2\beta c+c^2)}{\cos\tau}\,\frac\eta{1-\eta},\qquad
 B_\sigma=\frac{2e^{c^2/2}(6+12\beta c+8c^2+4\beta c^3+c^4)}{\cos\tau}\,\frac\eta{1-\eta}.
\]
Then for every phase and every \(a>0\)
\[
 |\sigma^2I_\mu(\vartheta)-J_\mu|\leq B_\mu,\qquad
 |\sigma^2I_\sigma(\vartheta)-J_\sigma|\leq B_\sigma,\qquad
 |\sigma^2I_{\mu\sigma}(\vartheta)|\leq\tfrac12(B_\mu+B_\sigma).
\]
With \(\tau=\pi/3\), for \(0<a\leq1\) both bounds are at most a polynomial
in \(1/a\) times \(e^{-11\pi^2/(18a^2)}\). Consequently every expansion of
\cref{thm:p15-fg-expansion}, including its eighth-order terms and the
expansions of \(J_\mu^{-1}\) and \(2J_\sigma^{-1}\), holds with
\(\sigma^2I_\mu(\vartheta)\) and \(\sigma^2I_\sigma(\vartheta)\) in place of
\(J_\mu\) and \(J_\sigma\), at every fixed phase, with a remainder bounded by
a constant times \(a^{10}\) uniformly in phase for all sufficiently small
\(a\); and the cross information tends to zero faster than every power of
\(a\).
\end{hbtheorem}

The theorem settles the question left by
\cref{thm:p15-w02b-p15-057-grouped-phase}: to every polynomial order the
information at a fixed grid equals its phase average, so Fisher's
small-width calculation of the information, as corrected in
\cref{sec:p15-fisher-grouping}, applies to the grid actually used and not
only to a randomized one. The bound is explicit but crude at moderate
widths. With \(\tau=\pi/4\) it gives \(0.080\) for location at \(a=1\) and
\(0.86\) at \(a=1.5\); direct summation of the cell series at \(20\)
phases gives periodic parts of about \(10^{-9}\) and \(8\times10^{-5}\).
These summations are illustrations; only the bound is proved. The
statement is about the information of the specified normal experiment; it
is not a statement about a sampling distribution at a fixed \(n\).

The third question concerns Fisher's comparison of the method of the
optimum with minimum \(\chi^2\). With \(\bar n_s=p_sN\) he puts
\(L=S(n_s\log(n_s/\bar n_s))\), minimized by the method of the optimum,
derives the equations (6) and (7) of the two methods, and shows that
\(L=\tfrac12\chi^2\) as a first approximation when every \(n_s-\bar n_s\)
is small compared with \(\bar n_s\); he then states that when \(\chi^2\)
closely approximates \(2L\) it follows a Type III curve
\(df\propto(\chi^2/2)^{(n'-3)/2}e^{-\chi^2/2}d(\chi^2/2)\), with \(n'\)
one more than the number of degrees of freedom
\cite[pp.~357--358]{Fisher1922Foundations}. His \(L\) is the \(D\) of
\cref{def:p15-fg-cells}, and Pearson's \(\chi^2\) is \(Q\) there. Two
printed formulas need correction. The denominator of \(\chi^2\) printed
on p.~357 is \(n_s\); Pearson's discrepancy, the one for which the next
display holds, has \(\bar n_s\). And that display, \(1+\chi^2=S(n_s^2/\bar
n_s)\), holds only with the counts expressed as proportions; with counts,
\(S(n_s^2/\bar n_s)=N+\chi^2\) because \(S(n_s)=S(\bar n_s)=N\). Equation
(7) is unaffected, since \(N\) does not depend on \(\theta\). The
deterministic comparison of \(2D\) and \(Q\) and the two score equations
are \cref{thm:p15-fg-deviance}. Under sampling, with fitted cells, the
comparison becomes a limit statement.

\begin{hbtheorem}{thm:p15-w02b-p15-057-chisq-sampling}{proved=annex:p15-w02b-p15-057.chisq-sampling}{Twice Fisher's \(L\) and Pearson's \(\chi^2\) under sampling with fitted cells}
Let \(K\geq2\) cells have true probabilities \(\pi_1,\ldots,\pi_K>0\), let
\((n_1,\ldots,n_K)\) be the counts of \(N\) independent observations, and
let \(\hat p_{N,1},\ldots,\hat p_{N,K}\) be random fitted cell
probabilities, functions of the counts, with \(\sum_s\hat p_{N,s}=1\) and
\(\sqrt N(\hat p_{N,s}-\pi_s)\) tight for each \(s\). Put
\(m_s=N\hat p_{N,s}\) and form \(D\) and \(Q\) of \cref{def:p15-fg-cells}
with these \(m_s\) on the event that every \(\hat p_{N,s}>0\). Then the
probability of that event tends to one, \(Q\) is tight, and
\(2D-Q\to0\) in probability as \(N\to\infty\). The same holds with
\(m_s=N\pi_s\).
\end{hbtheorem}

The tightness premise holds, for example, for the grouped normal estimate
of \cref{thm:p15-w02b-p15-057-grouped-mle}: each coordinate of
\(\sqrt N(\widehat\theta_N-\theta_0)\) converges in law by item~3 there and
so is tight, and the mean value theorem applied to the continuously
differentiable cell probabilities on a ball about \(\theta_0\), which
contains \(\widehat\theta_N\) with probability tending to one, bounds
\(\sqrt N|p_j(\widehat\theta_N)-p_j(\theta_0)|\) by a constant times
\(\sqrt N|\widehat\theta_N-\theta_0|\). The theorem justifies using \(\chi^2\) and \(2L\)
interchangeably in large samples with all expected counts growing in
proportion to \(N\); it gives nothing for sparse or empty cells, where
\cref{thm:p15-fg-deviance} already shows that the relative-error premise
fails. The Type III sampling law Fisher states for \(\chi^2\) needs more:
a joint limit for the whole vector of standardized cell counts, and the
projection that fitting induces on it. Both are supplied next, for an
estimate that behaves to first order like the maximum likelihood estimate
from the counts.

\begin{hblemma}{thm:p15-w02b-p15-057-lindeberg-vector}{proved=annex:p15-w02b-p15-057.lindeberg-vector}{A vector central limit statement for smooth functions}
Let \(\xi_1,\xi_2,\ldots\) be independent identically distributed random
vectors in \(\mathbb R^r\), \(r\geq1\), with \(|\xi_1|\leq b<\infty\),
\(E\xi_1=0\) and \(E\xi_1\xi_1^{\top}=I_r\), and put
\(V_N=N^{-1/2}\sum_{i\leq N}\xi_i\). Let \(G\) have independent standard
normal coordinates. If \(h:\mathbb R^r\to\mathbb R\) is three times
continuously differentiable, bounded, and has bounded second and third
partial derivatives, then \(Eh(V_N)\to Eh(G)\). Consequently
\(P(|V_N|^2\leq x)\to F_r(x)\) for every real \(x\), where \(F_r\) is the
distribution function of \(\chi^2_r\), the law
\(\mathrm{Gamma}(r/2,1/2)\) of \cref{thm:p15-student-orthogonal}.
\end{hblemma}

\begin{hbtheorem}{thm:p15-w02b-p15-057-chisq-law}{proved=annex:p15-w02b-p15-057.chisq-law}{The limiting law of \(\chi^2\) and \(2L\) with fitted cells}
Let \(K\geq2\) cells have probabilities \(p_1(\theta),\ldots,p_K(\theta)\),
\(\sum_sp_s(\theta)=1\), on an open set \(\Theta\subseteq\mathbb R^q\),
\(q\geq0\) (with \(q=0\) meaning that nothing is fitted). Let the counts
\((n_1,\ldots,n_K)\) come from \(N\) independent observations with cell
probabilities \(\pi_s=p_s(\theta_0)>0\), and suppose that \(p\) is
continuously differentiable near \(\theta_0\) and that the \(K\times q\)
matrix \(A\) with rows \(\pi_s^{-1/2}\nabla p_s(\theta_0)^{\top}\) has rank
\(q\). Put \(Z_{N,s}=(n_s-N\pi_s)/\sqrt{N\pi_s}\) and \(r=K-1-q\). Let
\(\widehat\theta_N\) be random vectors, functions of the counts, with
\[
 \sqrt N\,(\widehat\theta_N-\theta_0)-(A^{\top}A)^{-1}A^{\top}Z_N\to0
 \quad\text{in probability},
\]
and form \(Q\) and \(D\) of \cref{def:p15-fg-cells} with
\(m_s=Np_s(\widehat\theta_N)\) on the event that every
\(p_s(\widehat\theta_N)>0\), whose probability tends to one.
\begin{enumerate}
\item \(r\geq0\). If \(r\geq1\), then \(P(Q\leq x)\to F_r(x)\) and
 \(P(2D\leq x)\to F_r(x)\) for every real \(x\). If \(r=0\), then \(Q\to0\)
 and \(2D\to0\) in probability.
\item With \(q=0\) and \(m_s=N\pi_s\) this holds with \(r=K-1\).
\item The grouped normal estimate of
 \cref{thm:p15-w02b-p15-057-grouped-mle} (any maximizer, on the event that
 one exists) satisfies the premise with \(q=2\), so its \(\chi^2\) and
 \(2L\) have the limit law \(\chi^2_{K-3}\) when \(K\geq4\).
\end{enumerate}
\end{hbtheorem}

The law \(\mathrm{Gamma}(r/2,1/2)\) has density proportional to
\(x^{r/2-1}e^{-x/2}\); as a law for \(x/2\) it is Fisher's Type III curve
with exponent \((n'-3)/2\), where \(n'=r+1\)
\cite[p.~358]{Fisher1922Foundations}. Fisher's \(n'\) is one more than the
number of degrees of freedom in which the sample may differ from
expectation; the theorem identifies that number exactly as
\(r=K-1-q\): the counts lose one degree of freedom to their fixed total and
one to each fitted parameter, because the standardized residuals lie, to
first order, in the subspace orthogonal both to \((\sqrt{\pi_s})_s\) and to
the columns of \(A\). The law is a limit, not an exact law at any \(N\): the
statistic \(\chi^2\) has a discrete distribution, as Fisher notes for both
\(\chi^2\) and \(L\) \cite[p.~358]{Fisher1922Foundations}. It also needs an
estimate equivalent to first order to the maximum likelihood estimate from
the counts. An estimate obtained from the ungrouped observations, such as
the sample mean and standard deviation of the unrecorded values, does not
satisfy the premise, and the theorem gives no law for \(\chi^2\) fitted in
that way.

\section{Two discontinuous distributions: the Poisson series and the dilution method}
\label{sec:p15-w02b-p15-057-discrete}

Fisher ends his illustrations with discontinuous distributions. The first
is the Poisson series \cite[p.~359]{Fisher1922Foundations}; the third is the
dilution method for estimating the density of micro-organisms from the
presence or absence of growth on plates
\cite[pp.~363--366]{Fisher1922Foundations}. The grouped normal case between
them is \cref{sec:p15-w02b-p15-057-grouped}.

\begin{hbtheorem}{thm:p15-w02b-p15-057-poisson}{proved=annex:p15-w02b-p15-057.poisson}{The Poisson series}
Let \(X_1,\ldots,X_n\) be independent with
\(P_m(X=x)=e^{-m}m^x/x!\), \(x=0,1,2,\ldots\), and \(m\in(0,\infty)\).
\begin{enumerate}
\item If \(\overline X=n^{-1}\sum_iX_i>0\), the likelihood has the unique
 maximizer \(\widehat m=\overline X\), the unique root of
 \(\sum_i(X_i/m-1)=0\). If all \(X_i=0\), the likelihood increases as
 \(m\downarrow0\) and has no maximizer in \((0,\infty)\).
\item \(\sum_iX_i\) is sufficient for \(m\).
\item The family is regular at every \(m\), with \(I(m)=1/m\). The
 statistic \(\overline X\) is unbiased with
 \(\operatorname{Var}_m\overline X=m/n=1/(nI(m))\) for every \(n\), so its
 finite-sample information efficiency is one, and
 \(\sqrt n\,(\overline X-m)\Rightarrow N(0,m)\).
\end{enumerate}
\end{hbtheorem}

Item~3 is Fisher's \(\sigma_{\hat m}^2=m/n\)
\cite[p.~359]{Fisher1922Foundations}, which he obtains from the second
derivative of the log likelihood. It holds exactly at every sample size,
not only in the limit, because \(\overline X\) attains the bound of
\cref{thm:p15-w02b-p15-057-information}. The sample in which every count is
zero has positive probability \(e^{-nm}\) at every \(n\), and there the
formula \(\widehat m=\overline X\) gives the value \(0\), outside the
parameter set.

\begin{hbdefinition}{def:p15-w02b-p15-057-dilution}{The dilution experiment}
Fix a dilution factor \(a>1\), a finite nonempty set \(D\) of integers (the
dilutions used) and a number \(s\geq1\) of plates per dilution. The
parameter is the density \(n>0\) of organisms in the undiluted material,
written \(\theta=\log n\). At dilution \(x\) the expected number of
organisms on a plate is \(m_x=n\,a^{-x}\), and each plate is sterile with
probability \(p_x=e^{-m_x}\), fertile with probability \(q_x=1-p_x\), all
plates independently, as in Fisher's
\(\log p=-n/a^x\) \cite[p.~364]{Fisher1922Foundations}. The observation is
the number \(S_x\) of sterile plates at each \(x\in D\). Put
\[
 w(p)=\frac{p}{q}(\log p)^2,\qquad
 \mathcal W_D=\sum_{x\in D}w(p_x),\quad
 \mathcal D_D=-\sum_{x\in D}p_x\log p_x,\quad
 \mathcal V_D=\sum_{x\in D}p_xq_x,
\]
and write \(\mathcal W,\mathcal D,\mathcal V\) for the same sums over all
\(x\in\mathbb Z\). The \emph{sterile-count estimate} \(\widetilde\theta\)
is the solution of \(s\sum_{x\in D}p_x(\widetilde\theta)=T\), where
\(T=\sum_{x\in D}S_x\) is the total number of sterile plates.
\end{hbdefinition}

\begin{hbtheorem}{thm:p15-w02b-p15-057-dilution}{proved=annex:p15-w02b-p15-057.dilution}{Maximum likelihood and the sterile count in a dilution series}
In the experiment of \cref{def:p15-w02b-p15-057-dilution}, with \(D\)
fixed and \(s\to\infty\):
\begin{enumerate}
\item The information about \(\theta=\log n\) carried by one plate at
 dilution \(x\) is \(w(p_x)\), and that of the whole series is
 \(s\mathcal W_D\).
\item The log likelihood is strictly concave in \(\theta\). If at least one
 plate is sterile and at least one fertile, it has a unique maximizer
 \(\widehat\theta\), the unique root of Fisher's equation
 \(\sum_xS_x\log p_x-\sum_x(s-S_x)(p_x/q_x)\log p_x=0\); otherwise it
 has none. The first case has probability tending to one,
 \(\widehat\theta\to\theta\) almost surely, and
 \(\sqrt s\,(\widehat\theta-\theta)\Rightarrow N(0,1/\mathcal W_D)\).
\item If \(0<T<s|D|\), the sterile-count estimate exists and is unique;
 this has probability tending to one, and
 \(\sqrt s\,(\widetilde\theta-\theta)\Rightarrow
 N\bigl(0,\mathcal V_D/\mathcal D_D^2\bigr)\).
\item The efficiency of the sterile-count estimate relative to
 \(\widehat\theta\) is
 \(E_D(\theta)=\mathcal D_D^2/(\mathcal V_D\mathcal W_D)\leq1\), with
 equality only if \(p_x\) is the same at every \(x\in D\), that is, only
 if \(D\) has one element.
\end{enumerate}
\end{hbtheorem}

\begin{hbtheorem}{thm:p15-w02b-p15-057-dilution-average}{proved=annex:p15-w02b-p15-057.dilution-average}{The infinite series and the figure 87.71 per cent}
For every \(\theta\) the series \(\mathcal W,\mathcal D,\mathcal V\)
converge; as functions of \(\theta\) they are periodic with period
\(\log a\), and their averages over one period are
\[
 \overline{\mathcal W}=\frac{\pi^2}{6\log a},\qquad
 \overline{\mathcal D}=\frac1{\log a},\qquad
 \overline{\mathcal V}=\frac{\log2}{\log a}.
\]
When \(a=2\), \(\mathcal V(\theta)=1\) for every \(\theta\). For finite sets
\(D\) increasing to \(\mathbb Z\), \(E_D(\theta)\to\mathcal D^2/(\mathcal
V\mathcal W)\) at each \(\theta\). The ratio of the averages,
\[
 \frac{\overline{\mathcal D}^{\,2}}{\overline{\mathcal V}\,\overline{\mathcal W}}
 =\frac{6}{\pi^2\log2}=0.87705\ldots,
\]
does not depend on \(a\).
\end{hbtheorem}

These are the quantities of Fisher's argument. The weight \(w\) per plate
and the information \(s\sum w\) are his
\cite[p.~364]{Fisher1922Foundations}. The average \(\pi^2/(6\log a)\), its
value \(2.373138\) at \(a=2\), and the rule that one more sterile plate
lowers \(\log n\) by \(s^{-1}\log a\) are on p.~365 of
\cite{Fisher1922Foundations}; the rule is the reciprocal of the average
slope \(s\overline{\mathcal D}\) of the expected sterile count. The variance
\((\log a)^2s^{-1}\sum pq\) of the counting method, the value
\(\sum pq=1\) at \(a=2\) and the efficiency \(6/(\pi^2\log2)\), which he
gives as 87.71 per cent, are on p.~366 of \cite{Fisher1922Foundations}.

Two qualifications follow from the theorems. First, Fisher replaces
\(\mathcal W\) and \(\mathcal D\) by their averages over the phase
\(\theta\bmod\log a\) and divides averages; the efficiency at a given
density is \(\mathcal D^2/(\mathcal V\mathcal W)\), a periodic function of
\(\log n\), and the statement that the efficiency does not depend on the
dilution ratio holds for the ratio of averages, not for this function.
How far that function departs from the ratio of averages is bounded in
\cref{thm:p15-w02b-p15-057-dilution-oscillation} below, for every \(a\).
Second, the infinite series is an idealization: an infinite series
has infinitely many sterile plates, so the counting method needs a finite
\(D\), whose efficiency \(E_D\) approaches the infinite-series value only
as \(D\) grows. His statement that \(\sum pq\) is very nearly constant
is, at \(a=2\), exact. Recomputing his table for \(\log p=-2^{-x}\) (the
heading prints the exponent without its sign) reproduces the printed
entries to the printed precision except in the rows \(x=-4\), \(x=1\) and
\(x=4\); the row \(x=-4\) prints \(p=0.00000014167\), whereas
\(e^{-16}=1.1254\times10^{-7}\).

\begin{hbtheorem}{thm:p15-w02b-p15-057-dilution-oscillation}{proved=annex:p15-w02b-p15-057.dilution-oscillation}{The efficiency of the counting method at a given density}
In \cref{def:p15-w02b-p15-057-dilution} let \(\Lambda=\log a\),
\(\omega_k=2\pi k/\Lambda\), \(\gamma_k=|\Gamma(1-i\omega_k)|\), where
\(\Gamma(s)=\int_0^\infty t^{s-1}e^{-t}dt\) for \(\operatorname{Re}s>0\),
and \(\zeta(2)=\pi^2/6\).
\begin{enumerate}
\item For every \(\theta\),
 \[
 \begin{aligned}
  \Lambda\mathcal D(\theta)&=\sum_{k\in\mathbb Z}\Gamma(1-i\omega_k)
   e^{i\omega_k\theta},\qquad
  \Lambda\mathcal W(\theta)=\sum_{k\in\mathbb Z}\Gamma(2-i\omega_k)
   \zeta(2-i\omega_k)e^{i\omega_k\theta},\\
  \Lambda\mathcal V(\theta)&=\log2+\sum_{k\neq0}
   \frac{\Gamma(1-i\omega_k)(2^{i\omega_k}-1)}{i\omega_k}e^{i\omega_k\theta},
 \end{aligned}
 \]
 the series converging absolutely, with \(\zeta(s)=\sum_{j\geq1}j^{-s}\).
\item Put \(\delta_{\mathcal D}=2\sum_{k\geq1}\gamma_k\),
 \(\delta_{\mathcal W}=2\zeta(2)\sum_{k\geq1}(1+\omega_k^2)^{1/2}\gamma_k\)
 and \(\delta_{\mathcal V}=2\sum_{k\geq1}|2^{i\omega_k}-1|\gamma_k/\omega_k\).
 Then \(|\Lambda\mathcal D-1|\leq\delta_{\mathcal D}\),
 \(|\Lambda\mathcal W-\zeta(2)|\leq\delta_{\mathcal W}\) and
 \(|\Lambda\mathcal V-\log2|\leq\delta_{\mathcal V}\) for every \(\theta\).
 If \(\log2/\log a\) is an integer, that is \(a=2^{1/j}\) for a positive
 integer \(j\), then \(\delta_{\mathcal V}=0\) and \(\mathcal V\) is
 constant.
\item With \(b=e/\sin1\), \(|\Gamma(1-i\omega)|\leq
 b|\omega|e^{-\pi|\omega|/2}\) for \(|\omega|\geq1\). Hence, if
 \(a\leq e^{2\pi}\) (so that \(\omega_1\geq1\)), then with
 \(\rho=e^{-\pi^2/\Lambda}\),
 \[
  \delta_{\mathcal D}\leq\frac{2b\omega_1\rho}{(1-\rho)^2},\qquad
  \delta_{\mathcal W}\leq\frac{2\sqrt2\,\zeta(2)b\omega_1^2\rho(1+\rho)}{(1-\rho)^3},
  \qquad
  \delta_{\mathcal V}\leq\frac{4b\rho}{1-\rho}.
 \]
\item If \(\delta_{\mathcal D}<1\), \(\delta_{\mathcal V}<\log2\) and
 \(\delta_{\mathcal W}<\zeta(2)\), then for every density
 \[
  \frac{(1-\delta_{\mathcal D})^2}{(\log2+\delta_{\mathcal V})
  (\zeta(2)+\delta_{\mathcal W})}\leq
  \frac{\mathcal D^2}{\mathcal V\mathcal W}(\theta)\leq
  \min\Bigl\{1,\frac{(1+\delta_{\mathcal D})^2}{(\log2-\delta_{\mathcal V})
  (\zeta(2)-\delta_{\mathcal W})}\Bigr\}.
 \]
\item For \(a=2\): \(\delta_{\mathcal V}=0\),
 \(\delta_{\mathcal D}<3.9\times10^{-5}\),
 \(\delta_{\mathcal W}<8.1\times10^{-4}\), and the efficiency at every
 density lies in \([0.8765,0.8776]\). For \(a=3\):
 \(\delta_{\mathcal D}<4.7\times10^{-3}\),
 \(\delta_{\mathcal W}<6.2\times10^{-2}\),
 \(\delta_{\mathcal V}<1.7\times10^{-3}\), and the efficiency lies in
 \([0.835,0.923]\).
\end{enumerate}
\end{hbtheorem}

At \(a=2\) the efficiency at every density therefore lies between 87.65
and 87.76 per cent, so Fisher's 87.71 per cent \cite[p.~366]{Fisher1922Foundations}
holds at each density to within \(0.06\) percentage points, not only on
average. The bound of item~3 on \(|\Gamma(1-i\omega)|\) is cruder than the
exact value, by a factor of about \(4\) at \(a=2\); direct summation over
\(41\) phases (an illustration, not part of the proof) gives the narrower
range \(0.87699\) to \(0.87711\). For larger dilution factors the oscillation is genuine, and the
statement that the efficiency is independent of the dilution ratio holds
only for the ratio of averages. The bounds of item~3 decay like
\(e^{-\pi^2/\log a}\) and are informative only while that quantity is
small; at \(a=10\) item~4 gives no useful interval, and direct summation of
the series (an illustration, not a proof) shows efficiencies from about
\(0.78\) to \(0.97\) as the density varies over one period.

\section{Likelihood, confidence and posterior statements in one model}
\label{sec:p15-w02b-p15-057-three-statements}

The likelihood-estimation method ends with a maximizer and a variance
read from the curvature of the log likelihood. Such an interval is often
reported as if it carried a coverage guarantee or a posterior
probability. In the binomial model the three statements can be computed
exactly and compared. Throughout this section let \(X_1,\ldots,X_n\) be
independent Bernoulli variables with \(P_p(X_i=1)=p\), \(p\in[0,1]\), let
\(K=\sum_iX_i\), and let \(\Phi\) be the standard normal distribution
function.

\begin{hbdefinition}{def:p15-w02b-p15-057-wald-interval}{The curvature interval of the likelihood method}
The log likelihood is \(\ell_n(p)=K\log p+(n-K)\log(1-p)\) on \([0,1]\),
with \(0\log0=0\); its unique maximizer is \(\hat p=K/n\). The
per-observation information is \(I(p)=1/(p(1-p))\) for \(0<p<1\). The
\emph{expected-information variance} is \(v_E=\hat p(1-\hat p)/n\), which
is \(1/(nI(\hat p))\) when \(0<K<n\) and is \(0\) when \(K\in\{0,n\}\). The
\emph{observed-information variance} is \(v_O=-1/\ell_n''(\hat p)\), where
for \(K=0\) and \(K=n\) the second derivative is that of
\(n\log(1-p)\) and \(n\log p\) respectively, one-sided at the end point.
For \(z>0\) the curvature intervals are
\[
 W_E(K)=[\hat p-z\sqrt{v_E},\,\hat p+z\sqrt{v_E}]\cap[0,1],\qquad
 W_O(K)=[\hat p-z\sqrt{v_O},\,\hat p+z\sqrt{v_O}]\cap[0,1].
\]
Their \emph{coverage} at \(p\) is \(C(p)=P_p\{p\in W(K)\}\), the
fixed-parameter probability of \cref{def:p15-confidence-procedure}. For a
prior law on \([0,1]\), the \emph{posterior content} of \(W(k)\) is the
conditional probability of \(\{p\in W(k)\}\) given \(K=k\) in the joint
model of \cref{def:p15-bayes-experiment}.
\end{hbdefinition}

\begin{hbtheorem}{thm:p15-w02b-p15-057-three-statements}{proved=annex:p15-w02b-p15-057.three-statements}{The curvature interval is neither a coverage guarantee nor a posterior probability}
In the binomial model above, with \(n\geq1\) and \(z>0\):
\begin{enumerate}
\item \(v_E=v_O=\hat p(1-\hat p)/n\) when \(0<K<n\), whereas \(v_E=0\) and
 \(v_O=1/n\) when \(K\in\{0,n\}\).
\item The coverage of \(W_E\) is \(1\) at \(p=0\) and \(p=1\), is at most
 \(1-(1-p)^n-p^n\) for \(0<p<1\), and has infimum \(0\) over \(0<p<1\).
\item For each fixed \(p\in(0,1)\), the coverages of \(W_E\) and \(W_O\)
 both tend to \(2\Phi(z)-1\) as \(n\to\infty\).
\item Let \(n=4\) and \(\tfrac54\leq z<\sqrt3\). At \(p=5/8\) the coverage
 is \(1425/2048=0.6958\) for \(W_E\) and \(889/1024=0.8682\) for \(W_O\).
 At \(p=1/2\) it is \(7/8\) for \(W_E\) and \(1\) for \(W_O\).
\item Give \(p\) the uniform prior on \([0,1]\). The posterior law given
 \(K=k\) is \(\operatorname{beta}(k+1,n-k+1)\), and the posterior content
 of \(W_E(0)\) and of \(W_E(n)\) is \(0\). For \(n=4\) and \(z=z_{0.95}\),
 the number with \(\Phi(z_{0.95})=0.95\) (equal to \(1.6449\) to four
 decimals and so in the range of item~4), the posterior contents of
 \(W_E(1),W_E(2),W_E(3)\) are \(0.9176\), \(0.9878\), \(0.9176\), and those
 of \(W_O(0),W_O(4)\) are \(1-(1-z_{0.95}/2)^5=0.9998\), to four decimals.
\item For \(n=4\), \(K=2\) and \(0<z<2\), the interval \(W_E(2)=W_O(2)\) is
 exactly the set of \(p\) whose likelihood ratio \(L(p)/L(\hat p)\) is at
 least \((1-z^2/4)^2\).
\end{enumerate}
\end{hbtheorem}

\begin{hbexample}{def:p15-w02b-p15-057-three-example}{Four Bernoulli trials at nominal level 0.90}
Take \(n=4\) and \(z=z_{0.95}\), so that the large-sample coverage of
item~3 is \(0.90\). The \emph{likelihood statement} after \(K=2\) is: the
maximizer is \(1/2\), the curvature variance is \(1/16\), and the reported
interval \([0.0888,0.9112]\) is the set of values whose likelihood is at
least \((1-z_{0.95}^2/4)^2=0.1047\) times the maximum (item~6). This is a
comparison of values of \(p\) on the observed data; by itself it assigns
no probability to anything. The \emph{confidence statement} concerns the
procedure \(k\mapsto W(k)\) at a fixed \(p\): at \(p=5/8\) it covers with
probability \(0.6958\) for \(W_E\) and \(0.8682\) for \(W_O\), both below
\(0.90\), and the infimum for \(W_E\) is \(0\) (items~2 and~4). The
\emph{posterior statement} needs a prior; with the uniform prior the
reported intervals have posterior contents \(0\), \(0.9176\), \(0.9878\),
\(0.9176\), \(0\) for \(k=0,\ldots,4\) (item~5). None of the three numbers
is \(0.90\) in general, and they answer different questions. The exact
equal-tail interval of \cref{thm:p15-binomial-confidence-coverage} is the
construction that does guarantee coverage at least \(1-\alpha\) at every
\(p\), and \cref{thm:p15-confidence-credibility-boundary} shows that its
posterior content is a further, separate number.
\end{hbexample}

The large-sample agreement of item~3 is pointwise in \(p\), as
\cref{thm:p15-w02b-p15-057-asymptotic} is: it rests on \(0<p<1\), the
interior-point premise, and item~2 shows it is not uniform. The boundary
counts \(K\in\{0,n\}\), where the expected-information and
observed-information recipes disagree (item~1), are exactly the samples in
which the maximizer lies on the boundary of \([0,1]\), outside the
premises of that theorem.

\section{A misspecified family: what the maximizer and the curvature then estimate}
\label{sec:p15-w02b-p15-057-misspecified}

The theorems so far assume that the observations have law \(P_{\theta_0}\)
for some \(\theta_0\) in the model. When the sampling law \(Q\) is not in
the family, the likelihood can still be maximized, and the maximizer can
still converge, but to a value fixed by \(Q\) and the family jointly, and
the curvature of the log likelihood then estimates the wrong variance.
This section states that behaviour for one parameter and works it for
counts fitted by a Poisson family.

\begin{hbdefinition}{def:p15-w02b-p15-057-pseudo-true}{Misspecified sampling and the pseudo-true value}
In the setting of \cref{def:p15-w02b-p15-057-model}, let
\(X_1,X_2,\ldots\) be independent with a common law \(Q\) on
\((\mathsf X,\mathscr X)\). The family is \emph{misspecified for \(Q\)}
when \(Q\neq P_\theta\) for every \(\theta\in\Theta\). Write
\(\Lambda_Q(\theta)=E_Q\log f_\theta(X)\) whenever the positive part of
\(\log f_\theta(X)\) is \(Q\)-integrable, so that
\(\Lambda_Q(\theta)\in[-\infty,\infty)\). A \emph{pseudo-true value} is a
point \(\theta^*\in\Theta\) with \(\Lambda_Q(\theta^*)\) finite and
\(\Lambda_Q(\theta^*)=\sup_\theta\Lambda_Q(\theta)\). If \(Q\ll\mu\) with
density \(q\) and \(E_Q|\log q(X)|<\infty\), then
\(\Lambda_Q(\theta)=E_Q\log q(X)-E_Q\log\bigl(q(X)/f_\theta(X)\bigr)\)
wherever \(\Lambda_Q(\theta)\) is defined, so a pseudo-true value minimizes
the divergence of \(P_\theta\) from \(Q\) in the sense of
\cref{def:p15-w02b-p15-045-information}.
\end{hbdefinition}

\begin{hbtheorem}{thm:p15-w02b-p15-057-sandwich}{proved=annex:p15-w02b-p15-057.sandwich}{A root of the likelihood equation under a misspecified family}
Let \(\Theta\subseteq\mathbb R\) contain an open interval \(N\) about a
point \(\theta^*\), and let \(X_1,X_2,\ldots\) be independent with law
\(Q\). Assume:
\begin{enumerate}
\item[(M1)] there is a set \(A\) with \(Q(A)=1\) and \(f_\theta(x)>0\) for
 \(x\in A\), \(\theta\in N\);
\item[(M2)] for each \(x\in A\), \(\theta\mapsto\log f_\theta(x)\) is twice
 continuously differentiable on \(N\); write
 \(s_\theta(x)=\partial_\theta\log f_\theta(x)\);
\item[(M3)] \(E_Q|s_{\theta^*}(X)|<\infty\), \(E_Qs_{\theta^*}(X)=0\) and
 \(B=E_Qs_{\theta^*}(X)^2\in(0,\infty)\);
\item[(M4)] there are \(\delta>0\) and a measurable \(M\) with
 \(E_QM(X)<\infty\) and
 \(|\partial_\theta^2\log f_\theta(x)|\leq M(x)\) for \(x\in A\),
 \(|\theta-\theta^*|<\delta\); and
 \(J=-E_Q\partial_\theta^2\log f_{\theta^*}(X)>0\).
\end{enumerate}
Let \(\widehat\theta_n\) be random variables with
\(\widehat\theta_n\to\theta^*\) in probability and
\(Q(\ell_n'(\widehat\theta_n)=0)\to1\). Then:
\begin{enumerate}
\item \(\sqrt n\,(\widehat\theta_n-\theta^*)\Rightarrow N(0,B/J^2)\);
\item \(-\ell_n''(\widehat\theta_n)/n\to J\) in probability, so the
 curvature recipe \(n\cdot\bigl(-1/\ell_n''(\widehat\theta_n)\bigr)\)
 converges in probability to \(1/J\), which equals the limiting variance
 \(B/J^2\) if and only if \(B=J\);
\item if also \(E_QM(X)^2<\infty\), then
 \(\widehat B_n=n^{-1}\sum_{i\leq n}s_{\widehat\theta_n}(X_i)^2\to B\) in
 probability, and
 \(\widehat B_n\big/\bigl(\ell_n''(\widehat\theta_n)/n\bigr)^2\to B/J^2\)
 in probability;
\item if \(Q=P_{\theta^*}\) and the model is regular at \(\theta^*\)
 (\cref{def:p15-w02b-p15-057-regular}), then \(B=J=I(\theta^*)\) and
 items~1 and~2 are the conclusions of
 \cref{thm:p15-w02b-p15-057-asymptotic}.
\end{enumerate}
\end{hbtheorem}

Condition (M3), the vanishing of the mean score at \(\theta^*\), is the
first-order condition for an interior maximizer of \(\Lambda_Q\) when
differentiation under the expectation is allowed; it is stated as a
hypothesis because, with \(Q\) outside the family, the identity
\(\int\partial_\theta f_\theta\,d\mu=0\) of (R3) does not supply it. The
equality \(B=J\) on which the curvature recipe relies is the second
identity of (R3) read under the true law; under a misspecified family it
holds only by accident of \(Q\). The ratio \(B/J^2\) has the form of the
limit variance of the location estimates of
\cref{thm:p15-convex-location-normality}, with the likelihood score in
place of a location score function.

\begin{hbtheorem}{thm:p15-w02b-p15-057-poisson-misspecified}{proved=annex:p15-w02b-p15-057.poisson-misspecified}{Counts fitted by a Poisson family}
Let \(\mu\) be counting measure on \(\{0,1,2,\ldots\}\),
\(f_\lambda(x)=e^{-\lambda}\lambda^x/x!\) and \(\Theta=(0,\infty)\). Let
\(X_1,X_2,\ldots\) be independent with a law \(Q\) on
\(\{0,1,2,\ldots\}\) of mean \(\lambda^*\in(0,\infty)\) and variance
\(\sigma^2\in(0,\infty)\), not necessarily Poisson. Write
\(\overline X_n=n^{-1}\sum_{i\leq n}X_i\) and
\(S_n^2=n^{-1}\sum_{i\leq n}(X_i-\overline X_n)^2\).
\begin{enumerate}
\item \(\Lambda_Q(\lambda)=-\lambda+\lambda^*\log\lambda-E_Q\log X!\) is
 finite for every \(\lambda>0\), and \(\lambda^*\) is its unique
 maximizer, so \(\lambda^*\) is the unique pseudo-true value.
\item A maximum likelihood estimate exists if and only if
 \(\overline X_n>0\); it is then \(\overline X_n\), and it is unique.
 \(Q(\overline X_n>0)\to1\) and \(\overline X_n\to\lambda^*\) almost
 surely.
\item \(\sqrt n\,(\overline X_n-\lambda^*)\Rightarrow N(0,\sigma^2)\);
 hypotheses (M1)--(M4) hold at \(\theta^*=\lambda^*\), with
 \(E_QM(X)^2<\infty\), \(B=\sigma^2/\lambda^{*2}\) and
 \(J=1/\lambda^*\), so that \(B/J^2=\sigma^2\).
\item On \(\{\overline X_n>0\}\),
 \(-\ell_n''(\overline X_n)=n/\overline X_n\), so the curvature variance
 is \(\overline X_n/n\), and \(n\) times it tends to \(\lambda^*\) almost
 surely; the corrected estimate of item~3 of
 \cref{thm:p15-w02b-p15-057-sandwich} is exactly \(S_n^2\), which tends to
 \(\sigma^2\) almost surely.
\item For \(z>0\), the curvature interval
 \(\overline X_n\pm z\sqrt{\overline X_n/n}\) contains \(\lambda^*\) with
 probability tending to \(2\Phi\bigl(z\sqrt{\lambda^*}/\sigma\bigr)-1\),
 where \(\Phi\) is the standard normal distribution function; the
 interval \(\overline X_n\pm zS_n/\sqrt n\) contains \(\lambda^*\) with
 probability tending to \(2\Phi(z)-1\).
\item The fitted probability of a zero count, \(e^{-\overline X_n}\),
 tends almost surely to \(e^{-\lambda^*}\), which differs from \(Q(\{0\})\)
 unless \(Q(\{0\})=e^{-\lambda^*}\).
\end{enumerate}
\end{hbtheorem}

\begin{hbexample}{def:p15-w02b-p15-057-geometric-counts}{Geometric counts fitted as Poisson}
Let \(Q(\{k\})=2^{-k-1}\) for \(k=0,1,2,\ldots\). Then \(\lambda^*=1\) and
\(\sigma^2=2\) (computed in the annex proof of
\cref{thm:p15-w02b-p15-057-poisson-misspecified}). The Poisson maximizer
\(\overline X_n\) is a consistent estimate of the mean \(1\), with
limiting variance \(2/n\), but the curvature of the Poisson log likelihood
reports \(1/n\). At \(z=1.96\), whose nominal two-sided level is
\(2\Phi(1.96)-1=0.9500\) to four decimals, the curvature interval has
limiting coverage \(2\Phi(1.96/\sqrt2)-1=0.8342\) to four decimals, while
the interval with \(S_n\) has limiting coverage \(0.9500\). The fitted
chance of a zero count tends to \(e^{-1}=0.3679\) to four decimals,
whereas \(Q(\{0\})=1/2\). The maximizer estimates the mean correctly only
because the pseudo-true value of the Poisson family is the mean; the same
fit misstates every other feature of \(Q\) that the family ties to the
mean.
\end{hbexample}

The example separates three things that coincide under a correct model.
The maximizer has a limit (\(\lambda^*\)); that limit has a useful
interpretation (the mean of \(Q\)) only because of the particular family;
and the variance the likelihood reports (\(1/J\)) differs from the
sampling variance of the maximizer (\(B/J^2\)). The dependence
countermodel of \cref{thm:p15-fe-failures} misspecifies the sampling
scheme; the present case misspecifies the family and keeps the sampling
scheme correct.

\section{Limitations and what remains separate}
\label{sec:p15-w02b-p15-057-limits}

\begin{hbwarning}{def:p15-w02b-p15-057-limits}{Scope of the general theorems}
The asymptotic theorems are for independent identically distributed
observations. The scalar theorems and \cref{thm:p15-w02b-p15-057-multi-regular}
also need a parameter-free support; \cref{thm:p15-w02b-p15-057-multi}
allows a support depending on the parameter but only through its
conditions (S) and (C), which must be verified for each family, as
\cref{thm:p15-w02b-p15-057-typeiii-joint} does for the Type III curve with
\(p>1\). None of them covers the Type III regime \(-1<p\leq1\), where the
end point governs location and \cref{thm:p15-w02b-p15-057-endpoint} gives
the rate of the least observation instead, nor a root far from the true
value. They do
not cover dependent or nonidentically
distributed data, the dependence countermodel of
\cref{thm:p15-fe-failures} being one such case. Under a misspecified
family no parameter value gives the sampling law; a root of the likelihood
equation then converges, if at all, to a pseudo-true value, and the
curvature estimates the wrong variance, as
\cref{thm:p15-w02b-p15-057-sandwich} and
\cref{thm:p15-w02b-p15-057-poisson-misspecified} show. They give no confidence coverage or posterior probability; those
require \cref{def:p15-confidence-procedure} or
\cref{def:p15-bayes-experiment}, and
\cref{thm:p15-w02b-p15-057-three-statements} computes all three
statements in one model. Efficiency compares variances in a limit
and says nothing about a fixed-\(n\) error unless an exact calculation, as
in \cref{thm:p15-w02b-p15-057-typeiii} item~3, is available.
\end{hbwarning}

\index{likelihood!consistency theorem}
\index{likelihood!asymptotic normality}
\index{sufficiency!factorization criterion}
\index{efficiency!asymptotic variance}
\index{Fisher, R. A.!Type III efficiency}
\index{intrinsic accuracy}
